\begin{textsample}{POS Dim 2 – 2018\_c6 – Score 3.00 – 2018\_c6/enade\_2018\_tec\_des\_mod\_ce1d599a99.txt}  \label{ex:f2_pos_192}
CAPÍTULO 6 ANÁLISE TÉCNICA DA PROVA Este capítulo tem por objetivo apresentar o desempenho dos estudantes concluintes de Tecnologia em Design de Moda no Enade/2018.
Para isso, foram calculadas as estatísticas básicas da prova em seu todo ( seção 6.1.1 ), bem como as estatísticas relacionadas aos Componentes de Formação Geral ( seção 6.1.2 ) e de Conhecimento Específico da Área ( seção 6.1.3 ).
Dadas as suas características, foram analisadas, em separado, as questões objetivas ( seção 6.2 ) e as questões discursivas ( seção 6.3 ).
Para as questões objetivas, foram disponibilizados os índices de facilidade e de discriminação Ponto-Bisserial, também em separado, para os Componentes de Formação Geral ( seção 6.2.1 ) e de Conhecimento Específico ( seção 6.2.2 ).
De cada componente, uma das questões foi escolhida para exemplificar a análise gráfica, relacionando as alternativas escolhidas pelos estudantes ( inclusive o gabarito ) com o número de acertos no componente.
No Anexo I, apresenta -se a íntegra da análise gráfica para todas as questões objetivas.
Para cada uma das questões discursivas, os conteúdos dos tipos mais comuns de respostas dos estudantes são apresentados e comparados com o padrão de respostas esperado ( ver Anexo VIII com o padrão de respostas ).
Tomando -se como base as duas questões discursivas do Componente de Formação Geral, na seção 6.3.1.6 são apresentados comentários sobre a correção das respostas em relação à Língua Portuguesa.
Nas tabelas, constam as seguintes estatísticas das notas20 : média do desempenho na prova, erro padrão da média, desvio padrão, nota mínima, mediana e nota máxima para cada um de seus componentes.
Tais estatísticas contemplam o total de estudantes concluintes da Área de Tecnologia em Design de Moda inscritos e presentes à prova do Enade/2018, tendo em vista agregações, ou por Grandes Regiões e o país como um todo, ou por Categoria Administrativa, Organização Acadêmica da IES e Modalidade de Ensino.
Em relação aos gráficos de distribuição de notas, o intervalo considerado foi de 10 unidades, aberto à esquerda e fechado à direita, com exceção do primeiro intervalo, [ 0 ; 10 ], fechado em ambos os extremos.
Para os gráficos de distribuição das notas das questões discursivas, foram consideradas mais duas categorias : questão em branco21 e nota zero. 20 Uma definição dessas estatísticas pode ser encontrada no Glossário. 21 Nesse grupo estão incluídas também as respostas classificadas como nulas ou desconsideradas. 6.1 ESTATÍSTICAS BÁSICAS DA PROVA Esta seção apresenta estatísticas selecionadas e histogramas da nota geral ( 6.1.1 ) e de cada componente : Formação Geral ( 6.1.2 ) e Conhecimento Específico ( 6.1.3 ).
São também apresentadas estatísticas selecionadas de subpopulações, caracterizadas por Grande Região, Categoria Administrativa, Organização Acadêmica e Modalidade de Ensino. 6.1.1 Estatísticas Básicas Gerais Na Tabela 6.1, são apresentadas as Estatísticas Básicas da prova, por Grande Região, dos estudantes concluintes de Tecnologia em Design de Moda.
A Média das notas da prova, como um todo ( nas seções seguintes serão analisados os Componentes de Formação Geral e de Conhecimento Específico ) foi 45,8, sendo que os alunos da região Norte obtiveram a Média mais baixa ( 43,5 ), e os da região Sul obtiveram a Média mais alta ( 46,9 ).
As demais médias foram : 45,6 na região Nordeste, 44,3 na região Sudeste e 46,8 na região Centro-Oeste.
O Desvio padrão para o Brasil, como um todo, foi 13,2, sendo o maior Desvio padrão encontrado na região Centro-Oeste ( 14,0 ), e o menor, na região Norte ( 10,1 ), indicando uma dispersão um pouco menor das notas dessa última região.
A região que obteve a maior nota Máxima foi a Sul ( 86,9 ), ao passo que a região que atingiu a menor nota Máxima foi a Norte ( 60,9 ).
A Mediana do Brasil, como um todo, foi 45,8, sendo a maior Mediana obtida na região Centro-Oeste ( 47,6 ), e a menor, obtida na Norte ( 43,9 ).
A nota Mínima do Brasil ( 9,1 ) foi obtida na região Nordeste.
A Mínima na região Norte foi 25,1, a maior nota Mínima dentre as cinco Grandes Regiões.
Nas demais regiões as notas mínimas foram : 13,5 na Sudeste, 11,5 na Sul e 12,0 na Centro-Oeste.
Considerando -se as notas segundo Grande Região, observa -se que não existe diferença estatisticamente significativa ao nível de 95 \% entre as Médias obtidas pelas cinco Grandes Regiões²². ²² Todas as comparações deste capítulo utilizam os intervalos de 95 \%.
Os erros-padrão da média que possibilitam os testes estão disponíveis nas tabelas.
O comportamento das notas dos estudantes de todo o Brasil pode ser observado no Gráfico 6.1, que apresenta um histograma com a distribuição das mesmas : uma distribuição unimodal com a moda no intervalo ( 40 ; 50 ], seguido pelos intervalos ( 50 ; 60 ] e ( 30 ; 40 ].
Na Tabela 6.2, são apresentadas informações referentes à Média da nota final, desagregadas por Categoria Administrativa, por Organização Acadêmica e por Modalidade de Ensino.
A Média das notas da prova, como um todo, foi 45,8.
Em relação à Categoria Administrativa, os estudantes das IES Públicas obtiveram Média mais alta ( 48,2 ) que a Média nacional, e os das IES Privadas obtiveram Média mais baixa ( 45,4 ).
Observa -se que existe diferença estatisticamente significativa entre as médias das notas das IES Públicas e as das Privadas.
A diferença entre as médias das regiões Sul e Norte ( 3,4 ), a maior e a menor Média, é superior à diferença entre a das IES Públicas e a das Privadas ( 2,8 ), caracterizando -se uma maior diversidade regional do que administrativa.
O Desvio padrão das IES Públicas e das IES Privadas foi igual ao do Brasil, como um todo ( 13,2 ).
No tocante à Organização Acadêmica, os CEFET/IFET e as Universidades obtiveram Média mais alta que a nacional ( 48,7 e 46,3, respectivamente ).
A Média dos Centros Universitários e a das Faculdades foram menores do que a nacional ( 45,1 e 44,8, respectivamente ).
Constata -se que existe diferença estatisticamente significativa ao nível de 95 \% entre a Média mais alta, 48,7 dos CEFET/IFET, e a dos Centros Universitários e Faculdades.
A Média da Modalidade Educação a Distância ( 45,9 ) foi superior à Média da Educação Presencial ( 45,8 ), e a diferença entre elas não estatisticamente significativa, ao nível de 95 \%. 6.1.2 Estatísticas Básicas no Componente de Formação Geral Na Tabela 6.3, são apresentadas as Estatísticas Básicas em relação ao componente da prova que avalia a Formação Geral dos estudantes concluintes.
Os alunos de todo o Brasil obtiveram desempenho médio de 42,2.
Quanto à variabilidade, o Desvio padrão das notas dos estudantes do Brasil, como um todo, foi 17,8.
A maior Média foi obtida na região Sul ( 43,1 ), e a menor, na região Norte ( 38,2 ).
As demais médias foram : 41,4 na região Nordeste, 42,2 na região Sudeste e 41,0 na região Centro-Oeste.
Já o maior Desvio padrão foi obtido na região Centro-Oeste ( 19,9 ), e o menor, na região Norte ( 16,0 ).
Os demais desvios padrões foram : 16,4 na região Nordeste, 17,6 na região Sudeste e 18,6 na região Sul.
A maior nota no Componente de Formação Geral da prova do Enade/2018 foi 91,6, obtida por, pelo menos, um aluno na região Sul.
A menor nota Máxima foi obtida na região Norte ( 68,2 ).
A Mediana do Brasil, como um todo, foi 41,6, sendo a menor Mediana encontrada na região Norte ( 38,7 ), e a maior, na região Sul ( 42,6 ).
A nota Mínima nessa parte foi zero em quase todas as regiões, com exceção da região Norte, cuja nota Mínima foi 10,1.
Considerando -se as notas, segundo Grande Região, observa -se que não existe diferença estatisticamente significativa entre as Médias das notas do Componente de Formação Geral obtidas nas Cinco Grandes Regiões.
No Gráfico 6.2, é possível verificar a avaliação do desempenho dos estudantes do Componente de Formação Geral, a partir do histograma da distribuição das notas correspondentes.
A distribuição é unimodal, com moda em ( 30 ; 40 ], intervalo imediatamente anterior ao modal da distribuição de notas da prova, como um todo ( Gráfico 6.1 ).
Nota -se, ainda que, no Gráfico 6.2, as notas apresentam maior dispersão do que as no Gráfico 6.1 ( distribuição das notas da prova ), confirmado pela comparação dos desvios padrões : 13,2 para a nota da prova, como um todo, e 17,8 para o Componente de Formação Geral.
Norte ( 68,2 ).
A Mediana do Brasil, como um todo, foi 41,6, sendo a menor Mediana encontrada na região Norte ( 38,7 ), e a maior, na região Sul ( 42,6 ).
A nota Mínima nessa parte foi zero em quase todas as regiões, com exceção da região Norte, cuja nota Mínima foi 10,1.
Considerando -se as notas, segundo Grande Região, observa -se que não existe diferença estatisticamente significativa entre as Médias das notas do Componente de Formação Geral obtidas nas Cinco Grandes Regiões.
No Gráfico 6.2, é possível verificar a avaliação do desempenho dos estudantes do Componente de Formação Geral, a partir do histograma da distribuição das notas correspondentes.
A distribuição é unimodal, com moda em ( 30 ; 40 ], intervalo imediatamente anterior ao modal da distribuição de notas da prova, como um todo ( Gráfico 6.1 ).
Nota -se, ainda que, no Gráfico 6.2, as notas apresentam maior dispersão do que as no Gráfico 6.1 ( distribuição das notas da prova ), confirmado pela comparação dos desvios padrões : 13,2 para a nota da prova, como um todo, e 17,8 para o Componente de Formação Geral.
Na Tabela 6.4 são apresentadas as informações referentes ao desempenho dos concluintes do Componente de Formação Geral, em diferentes agregações : Categoria Administrativa, Organização Acadêmica e Modalidade de Ensino.
Observa -se que não existe diferença estatisticamente significativa entre as médias dos tipos de Categoria Administrativa : a maior Média obtida por estudantes de IES Públicas ( 44,0 ) e a menor pelos de IES Privadas ( 41,9 ).
Nota -se que também não há diferença estatisticamente significativa entre a maior das médias dos quatro tipos de Organização Acadêmica e a menor delas, o valor maior para os CEFET/IFET ( 45,0 ) e a menor para os Faculdades ( 40,9 ).
No que se refere à Modalidade de Ensino, a Educação a Distância apresentou maior Média do que a Educação Presencial ( médias 42,7 e 42,1, respectivamente ).
Tais médias não apresentam, tampouco, diferença estatisticamente significativa. 6.1.3 - Estatísticas Básicas do Componente de Conhecimento Específico Na Tabela 6.5, são apresentadas as Estatísticas Básicas referentes ao Componente de Conhecimento Específico da área de Tecnologia em Design de Moda.
A Média do desempenho dos alunos do Brasil, como um todo, foi 47,0.
A maior Média foi obtida na região Centro-Oeste ( 48,7 ), e a menor, na região Sudeste ( 45,0 ).
As demais médias foram : 45,2 na região Norte, 47,0 na região Nordeste e 48,1 na região Sul.
Quanto à variabilidade das notas, o Desvio padrão do Brasil, como um todo, foi 13,8, sendo o maior Desvio padrão observado nas regiões Sudeste e Centro-Oeste ( 14,2 ), e o menor, na região Norte ( 10,8 ).
Os demais desvios foram : 13,4 na região Nordeste e 13,7 na região Sul.
A Mediana das notas dos estudantes de todo o Brasil foi 47,2.
A maior Mediana ocorreu na região Centro-Oeste ( 49,9 ), e a menor, nas regiões Norte e Sudeste ( 45,2 ).
As demais medianas foram : 46,4 na região Nordeste e 48,5 na região Sul.
A nota Máxima do Brasil, como um todo, foi 87,7, sendo obtida por, pelo menos, um aluno na região Nordeste.
As demais notas máximas foram : 65,6 na região Norte, 81,3 na região Sudeste, 86,4 na região Sul e 83,8 na região Centro-Oeste.
A nota Mínima foi 7,1, obtida por, pelo menos, um aluno na região Nordeste.
As demais notas mínimas foram : 24,7 na região Norte, 9,6 na região Sudeste, 8,6 na região Sul e 14,2 na região Centro-Oeste.
Observa -se que não existe diferença estatisticamente significativa entre as duas menores Médias, das notas do Componente de Conhecimento Específico, das regiões Norte e Sudeste.
No entanto, existe diferença estatisticamente significativa entre as regiões Sudeste e Sul.
Assim como nos Gráficos 6.1 e 6.2, no Gráfico 6.3, apresentado a seguir, é possível ser feita uma avaliação do desempenho de concluintes de Tecnologia em Design de Moda em relação ao Componente de Conhecimento Específico, com um histograma da distribuição das notas correspondentes.
Esta também é uma distribuição unimodal, e o grupo modal é o ( 40 ; 50 ], o mesmo grupo modal da prova, como um todo, e acima do grupo modal para a Formação Geral.
Na Tabela 6.6, é apresentada uma comparação dos resultados em relação à Categoria Administrativa, à Organização Acadêmica e a Modalidade de Ensino, agora levando -se em conta o desempenho de alunos do Componente de Conhecimento Específico da prova.
No que se refere à Organização Acadêmica, a maior Média foi a dos CEFET/IFET ( 50,0 ), vindo a seguir a das Universidades ( 47,6 ), e, por fim, a dos Centros Universitários e das Faculdades ( 46,1 ).
O maior Desvio padrão, e acima do valor para o Brasil, como um todo, foi o das Universidades e o das Faculdades ( 14,0 ).
Os Centros Universitários obtiveram a maior nota Máxima ( 87,7 ).
As Universidades obtiveram nota Máxima 86,4 ; as Faculdades, 79,1 e os CEFET/IFET, 78,6.
As medianas foram : 48,5 nas Universidades, 46,4 nas Faculdades ; 45,7, nos Centros Universitários, a menor delas ; e 50,4 nos CEFET/IFET, a maior.
A nota Mínima foi 14,2 para os CEFET/IFET, 9,6 para as Faculdades, 8,6 para as Universidades e 7,1 para os Centros Universitários.
Observa -se que não existe diferença estatisticamente significativa ao nível de 95 \% no Componente de Conhecimento Específico entre as médias dos Centros Universitários e das Faculdades, as menores notas, mas existe entre estes e a dos CEFET/IFET, a maior Média.
Quanto à Categoria Administrativa, observa -se um comportamento semelhante àquele da prova, como um todo, ou seja, existe diferença estatisticamente significativa entre as médias das IES Públicas ( 49,5 ) e as das IES Privadas ( 46,5 ).
Nesse caso também, a maior Média foi obtida por alunos de IES Públicas de ensino.
Quanto ao comportamento das médias no que toca à Modalidade de Ensino, não há diferença estatisticamente significativa entre a Média da Educação Presencial ( 47,0 ) e a da Educação a Distância ( 46,9 ). 6.2 ANÁLISE DAS QUESTÕES OBJETIVAS Esta seção apresenta estatísticas selecionadas e histogramas das Questões Objetivas de Formação Geral ( 6.2.1 ) e de Conhecimento Específico ( 6.2.2 ).
São também apresentadas e comparadas as médias das subpopulações caracterizadas por Grande Região. 6.2.1 Componente de Formação Geral Na Tabela 6.7, são apresentadas as Estatísticas Básicas relativas às oito questões objetivas do componente da prova que abrange a Formação Geral dos estudantes.
A Média do Brasil foi 46,2.
A menor Média foi encontrada na região Norte ( 41,5 ), e a maior, na região Sul ( 47,0 ).
As demais médias foram : 46,5 na região Nordeste, 45,2 nas regiões Sudeste e Centro-Oeste.
O Desvio padrão do Brasil foi 23,0, sendo o maior Desvio padrão encontrado na região Centro-Oeste ( 25,7 ), e o menor, na região Nordeste ( 21,2 ).
Os demais desvios foram : 21,3 na região Norte, 23,4 na região Sudeste e 23,6 na região Sul.
As medianas do Brasil, como um todo, e de três das cinco regiões foram 50,0.
Nas regiões Norte e Centro-Oeste a Mediana foi 37,5.
A nota Máxima 100,0 foi alcançada em quatro das Grandes Regiões.
A exceção foi a região Norte, com nota Máxima 87,5.
As notas mínimas ( 0,0 ) foram iguais para quatro das cinco regiões, excetuando apenas a região Norte, com nota Mínima 12,5.
Na Tabela 6.8, são apresentados o Índice de Facilidade e o Índice de Discriminação ( Ponto-Bisserial ) para cada uma das questões objetivas do Componente de Formação Geral.
Quanto ao Índice de Facilidade, foram usadas as seguintes cores para diferenciar o nível de dificuldade da questão :  Azul para as questões classificadas com índice Muito fácil ( \&gt ; =0,86 ), verde para as questões classificadas com índice Fácil ( 0,61 a 0,85 ), amarelo para as questões classificadas com Médio ( 0,41 a 0,60 ), vermelho para as questões Já quanto ao Índice de Discriminação, foram usadas as seguintes cores para qualificar a questão :  As questões classificadas com índice Fraco receberam a cor vermelha ( \&lt ; =0,19 ), as classificadas com Médio receberam a cor amarela ( 0,20 a 0,29 ), as classificadas com Bom receberam a cor verde ( 0,30 a 0,39 ) e as classificadas com Muito bom ( \&gt ; =0,40 ) receberam a cor azul.
As questões objetivas do Componente de Formação Geral, segundo o Índice de Facilidade, foram assim avaliadas : das oito questões, nenhuma teve o Índice de Facilidade classificado como Muito fácil, e uma questão foi tida como Fácil, por ter índice de acertos 0,61.
Seis questões foram consideradas com índice de dificuldade Médio, situando -se no intervalo entre 0,42 e 0,57 do Índice de Facilidade, ou seja, houve entre 42,0 \% e 57,0 \% de acertos, enquanto uma questão foi classificada na categoria Difícil, com Índice de Facilidade 0,24.
Nenhuma questão apresentou menos de 15 \% de acertos, razão pela qual não houve questão classificada como Muito difícil.
O Índice de Facilidade variou de 0,24 a 0,61, e o de Discriminação, de 0,43 a 0,51.
Na Tabela 6.9, é apresentada a distribuição das questões, levando -se em conta simultaneamente a classificação dos dois índices.
As oito questões tiveram Índice de Discriminação Muito bom, e figuraram dentre três níveis de dificuldade Fácil, Médio e Difícil : uma classificada na categoria Fácil ( Questão 2 ) do Índice de Facilidade, seis na categoria Médio ( questões 1, 3, 5, 6, 7 e 8 ), e uma, na categoria Difícil ( Questão 4 ).
Em particular, a Questão 5 foi a que apresentou o maior poder discriminatório, com índice 0,51, e foi considerada Médio em termos de facilidade, com uma proporção de 0,45 acertos.
O máximo classificadas com Difícil ( 0,16 a 0,40 ) e roxo para as questões classificadas com Muito difícil ( \&lt ; =0,15 ).
Já quanto ao Índice de Discriminação, foram usadas as seguintes cores para qualificar a questão : • As questões classificadas com índice Fraco receberam a cor vermelha ( \&lt ; =0,19 ), as classificadas com Médio receberam a cor amarela ( 0,20 a 0,29 ), as classificadas com Bom receberam a cor verde ( 0,30 a 0,39 ) e as classificadas com Muito bom ( \&gt ; =0,40 ) receberam a cor azul.
As questões objetivas do Componente de Formação Geral, segundo o Índice de Facilidade, foram assim avaliadas : das oito questões, nenhuma teve o Índice de Facilidade classificado como Muito fácil, e uma questão foi tida como Fácil, por ter índice de acertos 0,61.
Seis questões foram consideradas com índice de dificuldade Médio, situando -se no intervalo entre 0,42 e 0,57 do Índice de Facilidade, ou seja, houve entre 42,0 \% e 57,0 \% de acertos, enquanto uma questão foi classificada na categoria Difícil, com Índice de Facilidade 0,24.
Nenhuma questão apresentou menos de 15 \% de acertos, razão pela qual não houve questão classificada como Muito difícil.
O Índice de Facilidade variou de 0,24 a 0,61, e o de Discriminação, de 0,43 a 0,51.
Na Tabela 6.9, é apresentada a distribuição das questões, levando -se em conta simultaneamente a classificação dos dois índices.
As oito questões tiveram Índice de Discriminação Muito bom, e figuraram dentre três níveis de dificuldade Fácil, Médio e Difícil : uma classificada na categoria Fácil ( Questão 2 ) do Índice de Facilidade, seis na categoria Médio ( questões 1, 3, 5, 6, 7 e 8 ), e uma, na categoria Difícil ( Questão 4 ).
Em particular, a Questão 5 foi a que apresentou o maior poder discriminatório, com índice 0,51, e foi considerada Médio em termos de facilidade, com uma proporção de 0,45 acertos.
O máximo de acertos foi alcançado pela Questão 2 com um Índice de Facilidade de 0,61 e índice de Discriminação 0,48.
A Questão 4 foi considerada Difícil, com índice de Facilidade 0,24 e Índice de Discriminação 0,46 ( Muito bom ). de acertos foi alcançado pela Questão 2 com um Índice de Facilidade de 0,61 e índice de Discriminação 0,48.
A Questão 4 foi considerada Difícil, com índice de Facilidade 0,24 e Índice de Discriminação 0,46 ( Muito bom ).
No Gráfico 6.4, para exemplificar, analisa -se o comportamento da questão de número 4 de Formação Geral.
Trata -se de uma questão considerada Difícil, em relação à facilidade e a que obteve Índice de Discriminação 0,46.
Neste gráfico, cada uma das cinco curvas representa o percentual de respostas em determinada alternativa da questão, em função do número de acertos dos estudantes nessa parte da prova ( Formação Geral/Múltipla Escolha ), antes de possíveis eliminações pelo critério do Ponto-Bisserial.
A curva em verde corresponde à alternativa B, a alternativa correta para essa questão.
Observa -se que dentre os estudantes com menor número de acertos, nessa parte do exame, a situação mais frequente foi a escolha da alternativa C ( em preto ), incorreta.
Por exemplo, dentre os estudantes que acertaram três questões, 33,6 \% escolheram a alternativa C ( em preto ), 21,4 \% escolheram a alternativa E ( em vermelho ), 19,2 \% escolheram a alternativa D ( em laranja ), 13,7 \% escolheram a alternativa A ( em azul ), e 11,4 \%, a B ( em verde ).
Entre os que acertaram três respostas entre as questões de múltipla escolha de Formação Geral, 0,4 \% deixou a questão em branco e 0,4 \% marcou mais de uma alternativa, invalidando a questão.
À medida que o número de acertos aumenta, indicando desempenho melhor nessa parte da prova, aumenta concomitantemente a proporção de estudantes que selecionaram a alternativa correta B, atingindo 100 \% para os estudantes com oito acertos.
Essa análise permite verificar como a questão discriminou os grupos de desempenho, justificando -se o alto índice obtido na questão.
Cumpre notar que não é possível inferir deste gráfico nem o Índice de Facilidade ( que seria uma média da proporção ponderada pela quantidade de alunos com cada uma das notas ), nem o Índice de Discriminação Ponto-Bisserial, por razão equivalente.
No caso extremo, no qual a grande concentração dos acertos dos alunos fosse abaixo de quatro, o Índice de Facilidade seria obrigatoriamente abaixo de 20 \% ( neste exemplo ).
Caso a concentração fosse em seis acertos ou mais, o índice seria obrigatoriamente acima de 40 \%.
Os gráficos relativos às demais questões de Formação Geral constam do Anexo I. 6.2.2 Componente de Conhecimento Específico Na Tabela 6.10, são apresentadas as Estatísticas Básicas em relação às questões objetivas do Componente de Conhecimento Específico da prova de Tecnologia em Design de Moda por Grande Região.
Nessa parte da prova, três questões objetivas do Componente de Conhecimentos Específicos não foram usadas no cômputo das notas por terem sido descartadas pelo critério do Ponto-Bisserial.
Assim, as notas foram calculadas com base em 24 das 27 questões objetivas de Conhecimentos Específicos.
A Média do Brasil deste componente foi 47,7.
A menor Média foi observada na região Sudeste ( 45,5 ), e a maior, na região Centro-Oeste ( 50,3 ).
O Desvio padrão de todo o Brasil foi 14,6, sendo o menor Desvio padrão encontrado na região Norte ( 11,9 ), e o maior, na região Sudeste ( 15,1 ).
A Mediana de todo o Brasil foi 45,8, o mesmo valor da Mediana encontrada nas regiões Norte, Nordeste e Sudeste.
Nas regiões Sul e Centro-Oeste, a Mediana foi 50,0.
A nota Máxima da prova foi 91,7, obtida nas questões objetivas do Componente de Conhecimento Específico, por, pelo menos, um aluno da região Nordeste.
Nas outras regiões as notas máximas foram : 66,7 na Norte, 83,3 nas regiões Sudeste e Centro-Oeste e 87,5 na região Sul.
A nota Mínima foi 4,2, obtida na região Sul.
Nas demais regiões, as notas mínimas foram : 25,0 na Norte, 16,7 na Centro-Oeste, e 8,3 nas regiões Nordeste e Sudeste.
Na Tabela 6.11, são apresentados os Índices de Facilidade e Discriminação ( Ponto-Bisserial ) das questões objetivas do Componente de Conhecimento Específico, para os estudantes de Tecnologia em Design de Moda.
Para facilitar a diferenciação das questões, serão usadas as mesmas cores da Tabela 6.8 para as diferentes classificações dos Índices de Facilidade e de Discriminação.
A partir do Índice de Facilidade obtido, observa -se que nenhuma das questões objetivas da prova foi classificada como Muito difícil, e quase metade ( 44,4 \% ) delas foram consideradas Difícil : 12, das 27 questões válidas.
Dez questões foram consideradas como Médio.
Uma questão foi classificada como Muito Fácil e quatro como Fácil.
Já quanto aos índices de discriminação das questões objetivas do Componente de Conhecimento Específico da prova, obtém -se como resultado a seguinte classificação : 15 das 27 questões foram consideradas boas, enquanto uma delas teve Índice de Discriminação Muito bom.
Assim, em quase três quintos das questões – 16 em 27 – os Índices de Discriminação foram Bom ou Muito bom.
Dentre as demais, oito delas foram classificadas como Médio, e outras três, como Fraco, sendo 11, por conseguinte, a quantidade de questões nos dois patamares mais baixos de discriminação.
Constata -se, assim, que a prova – no que se refere ao Componente de Conhecimento Específico – possuía média capacidade de discriminar dentre aqueles que dominam ou não o conteúdo.
Na Tabela 6.12, é apresentada a distribuição das questões levando -se em conta simultaneamente a classificação dos dois índices.
Dentre as questões que alcançaram os maiores índices de discriminação, uma delas foi classificada com Índice de Discriminação Muito bom, a de número 11.
Seu índice foi 0,42, e, quanto ao índice de Facilidade, foi classificada como Fácil, com 75 \% dos estudantes marcando a opção correta.
A Questão de número 23 foi a mais difícil dentre as 27 questões específicas válidas, com baixo Índice de Facilidade, apenas 21 \% de acertos.
Essa questão apresentou poder discriminatório muito baixo, 0,12, o que comprova ter sido a mais difícil para os estudantes.
Destaca -se, também, a Questão 26, com Índice de Facilidade 0,26, o que, em termos percentuais, corresponde a 26 \% de estudantes que responderam acertadamente.
Também foi 0,12 seu Índice de Discriminação.
Tais questões foram, portanto, pelo critério Ponto-Bisserial, consideradas inadequadas.
Por isso, as questões 23 e 26, além da Questão 12, mesmo sendo classificada como Fácil, 78 \% de acertos, foram eliminadas do cômputo da nota final.
A título de exemplo das análises do comportamento das questões objetivas, no Gráfico 6.5, analisa -se a Questão 11 do Componente de Conhecimento Específico.
Essa questão foi considerada, pelas respostas dos estudantes avaliados na prova, como uma questão com Índice de Facilidade Fácil, 0,75, ou seja, 75 \% dos estudantes assinalaram acertadamente a opção D, correspondente ao gabarito.
Como já comentado, seu Índice de Discriminação foi igual a 0,42, classificado como Muito bom, o maior valor de discriminação.
Neste gráfico, cada uma das cinco curvas representa o percentual de respostas em determinada alternativa da Questão 11, em função do número de acertos dos estudantes nessa parte da prova, antes de possíveis eliminações de questões pelo critério do Ponto-Bisserial.
A alternativa correta D, representada no gráfico pela curva em laranja, foi escolhida em maiores proporções pelos alunos com desempenho melhor nessa parte da prova.
Já as alternativas incorretas, também denominadas distratores, foram selecionadas, principalmente, por aqueles com notas mais baixas.
Observa -se que a soma não é 100 \%, por causa das questões não respondidas ou com mais de uma opção marcada.
Aqueles com nota zero ou um, na sua totalidade, deixaram esta questão em branco ou marcaram mais de uma alternativa, comportamento considerado inválido.
Observa -se que entre os participantes aqueles que optaram pela resposta correta D nesta questão foram só a partir de 3 acertos.
A proporção de alunos que selecionou a resposta correta D aumenta gradativamente, chegando a atingir 100 \% para 19, ou 20 acertos e depois com 22 ou mais acertos, enquanto a proporção dos que escolheram alternativas incorretas decai, a partir de três ou quatro acertos, em função do número de acertos nessa parte da prova.
Os gráficos relativos às demais questões do Conhecimento Específico constam do Anexo I. 6.3 ANÁLISE DAS QUESTÕES DISCURSIVAS Esta seção apresenta estatísticas selecionadas e histogramas das Questões Discursivas de Formação Geral ( 6.3.1 ) e Conhecimento Específico ( 6.3.2 ).
São também apresentadas e comparadas as médias de subpopulações, caracterizadas por Grande Região. 6.3.1 Componente de Formação Geral As análises dos resultados de desempenho dos estudantes de Tecnologia em Design de Moda nas duas questões discursivas relativas à Formação Geral encontram -se na Tabela 6.13 e no Gráfico 6.6.
Na Tabela 6.13, observa -se que a nota Média nesse conjunto de questões foi abaixo da obtida nas objetivas.
Os estudantes de todo Brasil obtiveram, em Formação Geral, Média 46,2 nas questões objetivas e 36,2 nas questões discursivas.
Pode -se notar também que o Desvio padrão nesse conjunto de questões foi menor do que o obtido nas objetivas : 23,0 nas questões objetivas e 19,9 nas questões discursivas.
A maior Média foi obtida na região Sudeste ( 37,7 ), e a menor, na região Norte ( 33,2 ).
A Média das notas das questões discursivas na região Nordeste foi 33,8, na região Sul igual a 37,2 e na região Centro-Oeste 34,7.
A Mediana de todo o Brasil, nesse componente, foi 38,0.
Nas regiões Norte ( 35,5 ), Nordeste e Centro-Oeste ( 36,0 em ambas ), os valores da Mediana foram menores que o valor encontrado para o Brasil.
Nas demais regiões, os valores da Mediana foram maiores quando comparados ao valor do Brasil : 38,5 na região Sudeste e 39,0 na região Sul.
A nota Máxima ( 91,5 ) foi obtida na região Sudeste, sendo 61,0 a Máxima na região Norte ; 84,0, na Nordeste ; 89,5, na Sul ; e 85,5, na região Centro-Oeste.
A nota Mínima ( 0,0 ) foi a mesma em todas as regiões do Brasil.
No Gráfico 6.6, está representada a distribuição das notas nas questões discursivas do Componente de Formação Geral.
A moda dessa distribuição ocorre no intervalo ( 30 ; 40 ], seguida pelo intervalo ( 40 ; 50 ].
Destaca -se também, como um máximo local, o intervalo [ 0 ; 10 ], com distribuição em torno de 13 \% do total de notas, sendo que no intervalo [ 0 ; 10 ] inclui -se, além da nota zero, a frequência de alunos que deixaram esse tipo de questão em branco.
Na sequência, os resultados verificados para cada uma das questões discursivas de Formação Geral serão apresentados, estabelecendo -se relações com os conteúdos abordados em cada uma delas.
Os comentários da Banca de docentes corretores, a respeito do observado na correção das respostas dos estudantes, suas impressões e conclusões serão apresentados junto à análise de cada questão.
Cumpre esclarecer que, tendo -se em vista que as questões discursivas de Formação Geral são padronizadas, ou seja, constam de todas as provas, os comentários da Banca são os mesmos para todas as carreiras acadêmicas, sendo direcionados a todos os estudantes que participaram do Enade/2018.
A seguir, serão analisados os desempenhos da Área de Tecnologia em Design de Moda nas duas questões discursivas de Formação Geral do Enade/2018, comparando -se os resultados obtidos com comentários para cada questão. 6.3.1.1 Análise de Conteúdo da Questão Discursiva 1 do Componente de Formação Geral Os dados de Tecnologia em Design de Moda, obtidos a partir das respostas à Questão 1, encontram -se na Tabela 6.14 e no Gráfico 6.7.
O desempenho médio dos estudantes na questão discursiva 1 ( média 30,3 ) foi inferior ao obtido na questão discursiva 2 ( média 33,5 ).
A maior Média para a Questão 1 foi obtida na região Centro-Oeste ( 31,5 ) e a menor na região Nordeste ( 28,1 ).
Quanto à variabilidade das notas, o Desvio padrão de todo o Brasil foi 25,1.
O menor Desvio padrão foi obtido na região Nordeste ( 24,3 ), e o maior Desvio padrão foi obtido na região Centro-Oeste ( 28,2 ).
A Mediana do Brasil como um todo foi 30,0, a mesma de quase todas as regiões, com exceção da região Norte ( 20,0 ).
A nota Máxima do Brasil ( 100,0 ) repetiu -se em quatro das cinco regiões, exceto na região Norte, com Máxima igual a 75,0.
Já a nota Mínima da questão discursiva 1 foi, sem exceção, a mesma para todas as regiões do Brasil : zero.
No Gráfico 6.7, mostra -se a distribuição das notas na questão discursiva 1 do Componente de Formação Geral.
Observa -se que a maior frequência corresponde ao intervalo ( 30 ; 40 ], com frequência igual a 16,9 \%, seguida pelo intervalo correspondente aos alunos que deixaram a resposta a essa questão em branco, 14,9 \% dos participantes.
Outros dois intervalos que se destacam são ( 20 ; 30 ] e ( 40 ; 50 ]. 6.3.1.2 Comentários sobre a correção de Conteúdo das respostas à Questão Discursiva 1 A Questão 1 exigia do aluno um posicionamento sobre conteúdo humanista e crítico.
O comando da questão solicitava a redação de um texto que formulasse e articulasse argumentos consistentes envolvendo temas previstos na Portaria nº 444, de 30 de maio de 2018, que dispõe sobre o componente de Formação geral do Enade/2018, tais como : solidariedade/violência, tolerância/intolerância, inclusão/exclusão, sexualidade, relações de gênero e relações étnico-raciais ( art. 7º, item VIII ), bem como ética, democracia e cidadania ( art. 7º, item I ).
O enunciado da questão buscava o estabelecimento da relação de três dimensões – os Direitos Humanos, a Democracia e as Fake News.
Para isso, recorreu a quatro suportes : a \textbf{reprodução} de uma foto e três fragmentos adaptados de textos que foram disponibilizados em páginas da Internet ( da Anistia Internacional, do Jornal El País-Brasil e do Labic-Laboratório de estudos sobre Imagem e Cibercultura ).
Os fragmentos traziam dados sobre o aumento do número de assassinatos de grupos minoritários, as origens e o trabalho da vereadora Marielle Franco, sua luta pelos Direitos Humanos, seu assassinato, e a produção de fake news.
O comando solicitava um texto abordando dois tópicos : o tensionamento entre a defesa dos Direitos Humanos e a produção de notícias falsas, tendo como base a vereadora Marielle – a defesa dos Direitos Humanos realizada por ela e as notícias falsas divulgadas após a sua morte.
As fake news também estavam presentes no segundo tópico de análise : os prejuízos que provocam para a sociedade democrática.
Por um lado, os textos selecionados para constituir o enunciado da questão continham apoio suficiente para formulação e articulação de argumentos para discorrer sobre o tensionamento entre a atuação da Marielle em defesa dos Direitos Humanos e fake news.
Por outro, havia ênfase nos aspectos biográficos e na trajetória da vereadora Marielle, o que levou a respostas pouco focadas no comando, que discutiam a vida e a morte da vereadora, sua importância como um ícone dos Direitos Humanos, opiniões em torno do seu assassinato, até mesmo sobre a responsabilidade pelo crime, muitos textos com a presença/defesa de fake news ou produção de novas.
O foco sobre Marielle também ensejou respostas basicamente centradas na questão da violência.
Foram frequentes as interpretações com viés político-partidário, não previstas no padrão de resposta.
O item do comando que solicitava a discussão sobre democracia e fake news encontrou menos apoio nos textos.
Apesar de se esperar que graduandos tivessem maturidade, inclusive política, para demonstrar alguma compreensão dos aspectos principais que envolvem o convívio democrático, esse aspecto foi pouco discutido e, quando presente, quase sempre foi apenas mencionado, numa quase \textbf{reprodução} do comando.
A falta de entendimento do que é o conceito de Direitos Humanos também foi percebida nas respostas, seja para recomendá -lo como importante para a sociedade, seja para criticá -lo como algo que defende bandidos.
Foi marcante a ideia que vincula Direitos Humanos a direitos apenas das minorias.
Não foi incomum encontrar respostas que personificavam os Direitos Humanos, tratando -os como sujeito de alguma ação que deveria ser implementada na sociedade em geral ou em casos específicos.
As fake news foram citadas muitas vezes, mas pouco explicadas, conforme previsto no padrão de resposta.
O argumento utilizado geralmente vinculava as consequências de fake news para a imagem da Marielle, ou seja, centrava no fato em si, já indicado no enunciado, sem qualquer análise ou aprofundamento próprio.
É importante salientar que muitas respostas, provavelmente de grupos vinculados a cursos específicos, continham argumentações recorrentes : os Direitos Humanos como algo vinculado à Constituição e a tratados dos quais o Brasil é signatário ; as fake news como produto do mundo digital e globalizado, relacionado à era da pós-verdade ; a discussão sobre o papel das mídias sociais e da imprensa.
Como já apontado, houve baixa quantidade de respostas contemplando o segundo tópico do comando – prejuízos da produção de notícias falsas para a sociedade democrática ; fato motivado, possivelmente, pela falta de um elemento que ressaltasse esse ponto nos textos motivadores.
De modo geral, os estudantes não conseguiram recorrer a conhecimentos tácitos sobre democracia para construir argumentos próprios.
As respostas que abordavam o tema, geralmente estavam vinculadas à ideia do direito à verdade, do direito à integridade da pessoa ou dos direitos gerais garantidos pela Constituição.
A relação da questão com o momento político brasileiro foi indicada em um número expressivo de respostas.
Por vezes, apenas questionando a temática da questão – caso Marielle – em uma prova para avaliar o ensino superior no Brasil ou a importância dos Direitos Humanos na atual conjuntura política do país, levando à anulação da resposta.
Houve também interpretações das consequências do uso das fake news no processo eleitoral de 2018, influenciando a escolha consciente e livre das pessoas.
As respostas demonstraram atenção político social frente a esses temas, porém com pouco domínio conceitual do campo de conhecimento sobre Direitos Humanos, fake news e democracia.
Em termos da linguagem adotada pelos alunos, há dois grandes grupos de respostas.
No primeiro grupo, os que responderam de forma clara, com autonomia de pensamento e explicitação de justificativas para os aspectos solicitados no enunciado, demonstrando conhecimento sobre o que significam Direitos Humanos, fake news e democracia, além da articulação entre essas questões.
Por outro lado, um segundo grupo dos que não conseguiram ir além do que foi citado no enunciado, apenas repetindo -o, muitas vezes desviando o foco da questão para temas correlatos, como a vida de Marielle ou a violência.
Nesse segundo grupo, mais facilmente foram encontradas respostas desconexas e mal articuladas, com textos confusos que explicitam o não conhecimento sobre o tema e a falta de adequação à linguagem, necessária a um concluinte do ensino superior.
As abordagens dos Direitos Humanos mais recorrentes nas respostas corretas foram : a referência da Marielle como representante/ícone dos DH ; e os DH como direitos previstos na Constituição e em tratados internacionais dos quais o Brasil é signatário.
No entanto, também foram encontradas críticas aos Direitos Humanos, tais como : DH vistos como “ direitos de bandidos ” ; a defesa pouco embasada de frases de efeito como “ direitos humanos para humanos direitos ”.
Também chamou atenção o tratamento dos DH como um personagem, ou um órgão público, evidenciando desconhecimento sobre o tema.
Em relação às fake news, dois argumentos previstos pelo padrão de respostas foram os mais presentes : a noção de que o mundo globalizado e a facilidade de \textbf{reprodução} de informações pelas mídias atuais levam à criação de mentiras que são compartilhadas e que geram consequências ; a Marielle como alvo de fake news que tentaram desconstruir sua luta e seu legado.
Também foi frequente a ideia de que a produção de informações falsas é realizada por uma imprensa tendenciosa com intenção de manipular a informação.
No item relacionado a fake news e democracia, as respostas corretas afirmavam que as fake news inviabilizam o acesso à verdade e à garantia da integridade de um Estado Democrático de Direito.
Também foi abordada a influência das fake news nos resultados das eleições dos últimos anos, nos Estados Unidos e no Brasil, gerando consequências para a vida das pessoas.
Por fim, um outro aspecto abordado nas respostas a este tópico foi a polarização da sociedade em termos políticos, com a afirmação de que as fake news dificultam o debate, estimulam posições extremas e rivalidades prejudiciais à convivência em uma democracia plena.
A notas mais fracas ( de zero a 35 ) foram cerca de 20 \% das respostas corrigidas – as que não estavam em branco nem foram desconsideradas ou anuladas.
Foram respostas que se restringiam a cópia ou paráfrases do enunciado ; as que tinham foco em desdobramentos políticos e policiais do caso Marielle ; as que tinham argumentação confusa e/ou incompleta.
A grande maioria das notas, em torno de 65 \%, foi mediana, entre 40 a 70.
Nesse grupo ficaram as respostas que continham alguma abordagem adequada sobre Direitos Humanos, sobre as atividades da vereadora Marielle e sobre as fake news divulgadas após sua morte.
As notas mais altas deste grupo foram para respostas que incluíram a questão do tensionamento, quase sempre, em um texto coerente e consistente.
Muitas vezes, a nota mediana foi consequência do baixo índice de respostas à segunda parte do comando, que solicitava a discussão sobre os prejuízos das fake news para a democracia.
As melhores notas só foram obtidas por aqueles que incluíram a discussão sobre os prejuízos das fake news para a democracia.
Apenas 15 \% das respostas corrigidas receberam avaliação igual ou superior a 75.
A deficiência principal detectada pelos corretores foi a dificuldade de grande parte dos alunos para interpretar o enunciado da questão e para se expressar em textos consistentes.
De modo geral, os temas foram tratados de forma fragmentada e pouco analítica.
Além disso, chamou à atenção a predominância do lugar comum na construção argumentativa das respostas, a superficialidade do debate em torno da democracia e o desconhecimento de conceitos como Direitos Humanos, esses confundidos com instituições e até pessoas. 6.3.1.3 Análise de Conteúdo da Questão Discursiva 2 do Componente de Formação Geral A Tabela 6.15 mostra que o desempenho médio dos estudantes na Questão discursiva 2 ( média 33,5 ) foi superior ao obtido na Questão discursiva 1 ( média 30,3 ).
A região Sul foi aquela cuja Média foi maior ( 34,8 ), enquanto a de menor Média foi a região Norte ( 27,4 ).
Quanto à variabilidade das notas, o Desvio padrão de todo o Brasil foi 22,5, inferior ao obtido na questão discursiva 1 ( 25,1 ).
O maior desvio nessa questão foi obtido na região Centro-Oeste ( 23,6 ), enquanto o menor foi obtido na região Norte ( 21,0 ).
A Mediana de todo o Brasil foi 35,0, a mesma de quatro regiões.
A exceção foi a região Norte, que obteve mediana 30,0.
A nota Máxima 100,0, foi obtida nas regiões Sudeste e Sul.
As demais notas máximas foram 70,0, na região Norte, e 85,0 nas regiões Nordeste e Centro-Oeste.
Já a nota Mínima da Questão discursiva 2 foi, sem exceção, a mesma para todas as regiões do Brasil : zero No Gráfico 6.8, mostra -se a distribuição das notas na Questão discursiva 2 do Componente de Formação Geral.
Observa -se que a maior frequência corresponde ao intervalo ( 40 ; 50 ], seguido de perto pelo intervalo ( 30 ; 40 ].
O terceiro intervalo de maior frequência foi o dos alunos que deixaram a questão em branco, com pouco mais de 13 \% dos participantes.
Nota -se, ainda, que as notas ficaram menos dispersas em comparação às da Questão discursiva de número 1, o que pode ser constatado também pela comparação do Desvio padrão das notas da Questão discursiva 2 ( 22,5 ) e o das notas da Questão discursiva 1 ( 25,1 ). 6.3.1.4 Comentários sobre a correção de Conteúdo das respostas à Questão Discursiva 2 Os textos e o comando da Questão 2 de Formação Geral são perfeitamente adequados ao que se espera avaliar em Formação Geral, ao tratar das funções sociais dos museus, contextualizando a temática sob o ponto de vista da memória e da perspectiva de futuro.
Na questão, estavam presentes, principalmente, o conteúdo humanista e crítico, a interpretação e a produção de textos sobre temas como cultura e arte, educação, ciência e democracia, previstos na Portaria que dispõe sobre o componente de Formação Geral do Enade/2018.
Os conhecimentos exigidos para responder ao comando, explicitados no padrão de resposta, vão além do senso comum, por se exigir a compreensão de conceitos como “ função social ” ( de instituições sociais ), “ memória ” ( no sentido de memória social ), “ perspectiva de futuro ” ( no contexto sócio-histórico-cultural ) e “ sociedade contemporânea ” ( que deveria independer da interpretação de cada grupo social a respeito da sociedade na qual vive ).
No entanto, apesar de envolver tais conceitos, espera -se que alunos concluintes de cursos superiores de graduação sejam capazes de discutir funções de instituições sociais, especialmente porque o enunciado trazia fato recente e de grande repercussão como apoio.
O padrão de resposta elencou as funções sociais que poderiam ser abordadas nas respostas dos estudantes, como também argumentos que poderiam ser usados ao longo das mesmas.
A tarefa dos estudantes era a construção de um texto expositivo que articulasse a importância de museus, suas funções sociais e as perspectivas de memória e de futuro, o que se mostrou de dificuldade de média a difícil.
Apesar de o enunciado ser claro, muitos estudantes, provavelmente influenciados pelos textos, limitaram -se a discutir aspectos relacionados ao Museu Nacional e, em particular, ao incêndio que o destruiu.
Foram poucos os que interpretaram corretamente o escopo do comando, que solicitava a discussão das funções sociais de museus de uma forma geral, no sentido de argumentar a importância deste tipo de instituição.
As quatro funções sociais elencadas no padrão de resposta foram encontradas na quase totalidade das redações dos estudantes de forma articulada com a importância dos museus.
Entretanto, quase nunca os textos abordavam mais do que uma função social e vinculavam, de forma clara, a importância desse tipo de instituição para a sociedade contemporânea.
Além disso, a solicitação de contextualização sob o foco museu/memória/perspectiva de futuro parece não ter sido bem compreendida e, em consequência, esteve pouco presente nas respostas.
Também se observou confusão entre os conceitos de história e memória e uma visão estreita do significado de sociedade contemporânea, às vezes, restrita à cidade do respondente, outras, ao Rio de Janeiro, ou ainda a um tipo de sociedade idealizada.
A relevância que os textos de apoio deram ao caso particular do incêndio do Museu Nacional do Rio de Janeiro e aos prejuízos causados ao seu acervo, induziu a maioria dos estudantes a elaborar respostas tratando das funções a partir da perda e, em particular, à perda do tipo de acervo daquele museu e do tipo de pesquisa que lá se desenvolve, divergindo do comando da questão, que solicitava texto abordando funções sociais de museus em geral.
Também foram encontradas muitas respostas que se limitavam a discutir o que o estudante considerava como possíveis causas do incêndio, a falta de manutenção, questões políticas e de gestão, sem responder ao que foi solicitado.
Foi também recorrente a transcrição de frases dos textos do enunciado nas respostas, diretamente ou por paráfrases, revelando certa dificuldade de criação textual.
A análise das respostas dos estudantes, além de revelar dificuldade de abordar a temática esperada, revelou outros dois aspectos predominantes.
Um deles, como já comentado, foi a ênfase no incêndio do Museu Nacional, que orientou grande parte dos textos dos estudantes.
Outro, também evidenciando desvio do esperado, foi o foco na perda patrimonial.
Nesse grupo, houve respostas que até apresentaram uma função social, mas olhada pelo ponto de vista da perda.
Nessa linha, muitas respostas enalteciam a importância do Museu Nacional, sua relevância no cenário nacional, o impacto da perda de seu acervo, isso complementado por frágeis reflexões sobre as funções sociais.
Em consequência, a perspectiva de desdobramentos mais ricos, que tratassem de aspectos sociais relacionados a diferentes tipos de museu foi prejudicada.
Muitas respostas trataram da preservação da memória e/ou guarda do acervo ; identidade, tanto individual quanto coletiva ; conhecimento sobre antepassados ; evolução da sociedade e dos seus costumes ; valorização e disseminação do conhecimento por meio da liberação do acesso aos acervos ; bem como atividades \textbf{pedagógicas} e exposições de arte.
O aspecto produção de pesquisa e conhecimento, embora frequente, foi mencionado de forma muito simples, sem qualquer aprofundamento, para além do que já era mencionado no enunciado.
Frequentemente apresentava -se o museu como produtor de pesquisas para o campo da saúde, o que evidencia desconhecimento das reais possibilidades de pesquisa, mesmo no Museu Nacional.
Quanto a perspectiva de futuro, quando presente no texto, era mencionada de forma breve, com afirmações genéricas sobre a possibilidade de melhoria do futuro pelo simples conhecimento do passado, numa percepção de que o passado determina o futuro, de forma linear.
Foram raras as respostas com alguma argumentação que associasse pesquisa, memória, identidade e reflexão social como base para a construção de alternativas de futuro.
A relação entre turismo e museu apareceu considerando a visita a museus como forma de lazer e, algumas vezes, como geradora de recursos.
O desempenho pode ser considerado fraco.
Muitas respostas deixavam claro que o estudante nunca tinha visitado um museu e, de uma forma geral, observou -se desconhecimento sobre suas funções sociais.
Não se esperava que história, memória social e museologia fossem assuntos dominados pelos estudantes das diversas áreas que participaram do Enade/2018, mas esperava -se que essa temática pudesse ser abordada do ponto de vista das experiências pessoais ou de conhecimento social.
Como uma vivência própria com os museus também se revelou frágil ou inexistente, o conhecimento sobre o assunto ficou bem restrito aos textos motivadores e ao senso comum, construído, basicamente, pelo noticiário em torno do incêndio.
O universo dos argumentos dos estudantes se mostrou muito limitado.
Cada dimensão foi praticamente representada por uma ideia simplista.
De forma recorrente : a preservação da memória para construção da identidade de um povo ; a compreensão do passado por meio das pesquisas como oportunidade para melhorar o futuro ; o museu como um espaço em que há um aprendizado estimulado pelo lazer.
Houve, também, uma grande concentração de respostas associando o museu a um local de guarda de acervo e de história.
As respostas à Questão 2 evidenciaram dificuldades conceituais.
Além dos conceitos de história, memória e museu se misturarem, parecendo não se distinguirem, ao mencionarem a pesquisa, também se observou falta de distinção entre pesquisas colegiais, com objetivos de aprendizado, e as realizadas por pesquisadores, com objetivos de avançar cientificamente.
A falta de compreensão do que se pedia no comando da questão e a dificuldade de expressão linguística foram dois aspectos marcantes na avaliação dos corretores.
Foi frequente os textos apresentarem ideias soltas ou de forma confusa, sem constituir unidade.
O vocabulário usado se revelou bem restrito, e fortemente apoiado no coloquial, na linguagem oral.
Predominou o senso comum em lugar de uma reflexão mais formal envolvendo os conceitos envolvidos. 6.3.1.5 Análise de Língua Portuguesa das Questões Discursivas do Componente de Formação Geral Os dados de Tecnologia em Design de Moda, obtidos a partir das respostas às questões discursivas do Componente de Formação Geral, no que tange à Língua Portuguesa, encontram -se na Tabela 6.16 e no Gráfico 6.9.
Nesse aspecto, os alunos de todo o Brasil, obtiveram Média 53,5.
A maior Média em relação à Língua Portuguesa foi obtida na região Sudeste ( 56,8 ), e a menor, na região Centro-Oeste ( 49,4 ).
Quanto à variabilidade das notas, o Desvio padrão de todo o Brasil foi 23,5.
O menor Desvio padrão foi obtido na região Sudeste ( 21,3 ) e o maior Desvio padrão foi obtido na região Centro-Oeste ( 26,2 ).
A Mediana das notas de Língua Portuguesa foi 62,5 para o Brasil, a mesma obtida nas regiões Norte, Sudeste e Sul.
A Mediana da região Nordeste foi 57,5, a menor, enquanto a da Centro-Oeste foi 60,0.
A nota Máxima para todo o Brasil foi 85,0, obtida por, pelo menos, um aluno na região Sul.
Na região Norte, a nota Máxima foi 77,5 ( a menor ), e, nas demais, 82,5.
Já a nota Mínima foi zero em todas as regiões do país. das pesquisas como oportunidade para melhorar o futuro ; o museu como um espaço em que há um aprendizado estimulado pelo lazer.
Houve, também, uma grande concentração de respostas associando o museu a um local de guarda de acervo e de história.
As respostas à Questão 2 evidenciaram dificuldades conceituais.
Além dos conceitos de história, memória e museu se misturarem, parecendo não se distinguirem, ao mencionarem a pesquisa, também se observou falta de distinção entre pesquisas colegiais, com objetivos de aprendizado, e as realizadas por pesquisadores, com objetivos de avançar cientificamente.
A falta de compreensão do que se pedia no comando da questão e a dificuldade de expressão linguística foram dois aspectos marcantes na avaliação dos corretores.
Foi frequente os textos apresentarem ideias soltas ou de forma confusa, sem constituir unidade.
O vocabulário usado se revelou bem restrito, e fortemente apoiado no coloquial, na linguagem oral.
Predominou o senso comum em lugar de uma reflexão mais formal envolvendo os conceitos envolvidos. 6.3.1.5 Análise de Língua Portuguesa das Questões Discursivas do Componente de Formação Geral Os dados de Tecnologia em Design de Moda, obtidos a partir das respostas às questões discursivas do Componente de Formação Geral, no que tange à Língua Portuguesa, encontram -se na Tabela 6.16 e no Gráfico 6.9.
Nesse aspecto, os alunos de todo o Brasil, obtiveram Média 53,5.
A maior Média em relação à Língua Portuguesa foi obtida na região Sudeste ( 56,8 ), e a menor, na região Centro-Oeste ( 49,4 ).
Quanto à variabilidade das notas, o Desvio padrão de todo o Brasil foi 23,5.
O menor Desvio padrão foi obtido na região Sudeste ( 21,3 ) e o maior Desvio padrão foi obtido na região Centro-Oeste ( 26,2 ).
A Mediana das notas de Língua Portuguesa foi 62,5 para o Brasil, a mesma obtida nas regiões Norte, Sudeste e Sul.
A Mediana da região Nordeste foi 57,5, a menor, enquanto a da Centro-Oeste foi 60,0.
A nota Máxima para todo o Brasil foi 85,0, obtida por, pelo menos, um aluno na região Sul.
Na região Norte, a nota Máxima foi 77,5 ( a menor ), e, nas demais, 82,5.
Já a nota Mínima foi zero em todas as regiões do país.
No Gráfico 6.9, mostra -se a distribuição das notas de Língua Portuguesa do Componente de Formação Geral.
Observa -se que a maior frequência ( 30,5 \% ) corresponde à dos alunos que obtiveram nota no intervalo ( 60 ; 70 ], seguido pelos intervalos ( 70 ; 80 ] e ( 50 ; 60 ].
Destacam -se, também, os alunos que deixaram ambas as questões em branco, representando quase 10 \% do total, caracterizando -se como máximo local. 6.3.1.6 Comentários sobre a correção das respostas de Formação Geral com respeito à Língua Portuguesa As duas questões discursivas de Formação Geral ensejaram duas configurações textuais distintas, tendo em vista os seus encaminhamentos : a Questão 1 teve um caráter mais opinativo e a Questão 2, mais expositivo, conforme se explicita a seguir.
Ambas as questões permitiram a observação do desempenho linguístico dos participantes em situação formal, competência exigida no processo de formação universitária.
As habilidades de produção textual envolvidas têm, igualmente, como condição básica, o domínio do padrão formal culto da Língua Portuguesa, no que diz respeito aos aspectos relativos às questões de caráter ortográfico, textual, morfossintático e vocabular.
O encaminhamento das questões evidencia a expectativa de que o participante utilizasse seus conhecimentos sobre os dois temas e estruturasse seus textos de acordo com as características do registro formal adequado à situação comunicativa.
Tal configuração determina exigências quanto : à adequação da seleção vocabular, ao desenvolvimento do conteúdo, à estruturação sintática dos períodos, à organização lógica das ideias, à utilização de procedimentos de encadeamento textual e de referenciação, à obediência às exigências morfossintáticas próprias da modalidade escrita da norma-padrão, ao respeito às regras ortográficas e às regras de acentuação gráfica.
Esse encaminhamento exigiu do graduando o domínio do modo de organização textual de base dissertativo-argumentativa, essencial ao processo de formação universitária.
As habilidades de produção textual envolvidas têm como condição básica o domínio do padrão formal culto da Língua Portuguesa, no que diz respeito aos aspectos relativos às questões de caráter ortográfico, textual, morfossintático e vocabular.
O enunciado da Questão 1 tomava como base três fragmentos de textos jornalísticos publicados em sites da internet, que abordavam o aumento do assassinato de minorias e de defensores de direitos humanos no Brasil, entre eles a vereadora Marielle Franco.
O comando solicitava que o participante desenvolvesse um texto dissertativo e argumentativo com as temáticas Fake News, Democracia e Direitos Humanos.
Embora não tenha havido referência explícita à tipologia textual esperada para o desenvolvimento da resposta, essa proposta encaminhou o participante para a elaboração de um texto opinativo, com o objetivo de discutir o problema e manifestar um posicionamento crítico.
O enunciado da Questão 2 abordava o incêndio do Museu Nacional do Rio de Janeiro, ocorrido em 2018, e seus impactos na ciência brasileira e internacional, tomando como base dois fragmentos de textos jornalísticos publicados em sites da internet.
O comando solicitava que o participante desenvolvesse um texto abordando três aspectos da função social dos museus na sociedade contemporânea, encaminhando -o para a elaboração de um texto expositivo que confirmasse a posição defendida nos fragmentos motivadores a respeito da importância dessas instituições.
Para atender a essas exigências, o padrão de resposta utilizado na avaliação das questões discursivas de Formação Geral considerou aspectos relevantes ao bom desempenho linguístico como competências distintas, de modo a permitir um mapeamento detalhado do domínio dos recursos disponíveis na Língua Portuguesa para a comunicação escrita formal : ortográficos, textuais, morfossintáticos e vocabulares ( ver detalhes no Anexo VIII – Padrão de Respostas ).
Aspectos Ortográficos Para avaliar a competência relativa ao domínio das convenções ortográficas da norma-padrão da Língua Portuguesa, observou -se o respeito às regras de acentuação gráfica e da grafia padrão das palavras ( com ausência de abreviaturas próprias da linguagem da internet ), de acordo com as convenções estabelecidas pela legislação em vigor e consubstanciadas no Vocabulário Ortográfico da Língua Portuguesa, editado pela Academia Brasileira de Letras.
Aspectos Textuais Para avaliar a competência relativa ao domínio dos procedimentos de estruturação textual, procurou -se observar a utilização de mecanismos linguístico-discursivos que estabeleçam o encadeamento lógico entre as partes do texto, de forma a garantir a progressão e a coerência textuais.
Esta competência destaca -se por sua relevância, já que são avaliadas a clareza e a unidade textual.
Aspectos Morfossintáticos Para avaliar o domínio dos diferentes aspectos morfossintáticos próprios da modalidade escrita formal da norma-padrão da Língua Portuguesa, observou -se se o texto produzido atendeu às seguintes exigências : a concordância nominal, a concordância verbal, a regência nominal, a regência verbal, a flexão nominal, a flexão verbal, a correlação entre os tempos verbais, a colocação pronominal e a utilização de sinais de pontuação que contribuam para a organização lógica da frase e do texto.
Aspectos Vocabulares Para avaliar a adequação da seleção vocabular à modalidade escrita formal da Língua Portuguesa, exigida pela situação comunicativa, observou -se o atendimento às seguintes exigências : precisão na seleção/utilização do vocabulário relacionado à temática solicitada pela questão ; ausência de marcas da oralidade, como termos de sentido muito genérico ( “ coisa ”, “ negócio ”, “ você ” ) e termos de registros mais informais ( como gírias, jargões, frases feitas, ditados populares, termos regionais ).
Assim, espera -se que o participante respeite a adequação vocabular não usando gírias ou expressões coloquiais, evite repetição desnecessária de palavras e utilize um vocabulário mais formal, como solicitado por um texto dissertativo.
Para efeito de pontuação, a avaliação do desempenho linguístico considerou três grandes grupos de competências, agrupando aspectos explicitados, com os seguintes pesos relativos : aspectos ortográficos ( 20 \% ) ; aspectos textuais ( 40 \% ), aspectos morfossintáticos e vocabulares ( 40 \% ).
Os critérios de pontuação utilizados no padrão de resposta procuraram atender à diversidade de desempenho dos graduandos.
Nos textos analisados, observaram -se, assim, vários níveis de atendimento às exigências do padrão formal próprio da modalidade escrita da Língua Portuguesa, configurando um"continuum".
Por um lado, textos que traduzem com exatidão a expectativa de uso formal da língua, em todos os níveis analisados, marcados por clareza e unidade ; e, no outro extremo, textos quase incompreensíveis devido ao grande número de desvios de todos os tipos, principalmente os de caráter textual.
As duas questões discursivas de Formação Geral analisadas revelaram desempenhos distintos dos participantes.
A Questão 1 propiciou a oportunidade de maior desenvolvimento do tema solicitado, gerando consequências na estruturação textual, já que os textos foram mais longos, mais elaborados e fluentes.
Em virtude da polêmica sobre a existência das fake news e sobre a visão pública da personagem da vereadora Marielle Franco, observaram -se respostas que polemizaram o tema, favorecendo a construção argumentativa do texto.
A Questão 2 teve um comportamento distinto, devido ao caráter mais expositivo do enunciado e ao reduzido conhecimento dos participantes sobre o tema.
Assim, os textos foram mais curtos, frequentemente com repetição de ideias dos textos motivadores.
Quanto aos aspectos linguísticos analisados durante esta avaliação, observaram -se os seguintes resultados : Aspectos ortográficos : O desempenho dos participantes revelou uma diferença muito grande nos dois aspectos analisados nesta competência : baixo índice de desvios da grafia padrão e grande índice de desvios de acentuação.
Em vários casos, ocorre ausência completa de acentuação gráfica em todas as palavras do texto.
Os resultados revelaram que a tendência dominante entre os universitários brasileiros é a eliminação da acentuação gráfica, talvez motivada pelos hábitos relacionados às redes sociais e pela ausência de esclarecimento dos meios de comunicação, das autoridades e das escolas sobre as decisões do Acordo Ortográfico da Língua Portuguesa de 1990.
Os casos mais sistemáticos de eliminação do acento indicador da sílaba tônica foram : - palavras proparoxítonas ( por exemplo “ generos ”, “ arbitro ”, “ politico ”, “ publicas ”, “ numero ” ) ; - palavras paroxítonas terminadas em ditongo crescente ( por exemplo “ area ”, “ noticiarios ”, “ individuo ”, “ historia ”, “ varios ”, “ noticia ”, “ negligencia ”, “ ocorrencia ” ) ; - palavras oxítonas ( por exemplo “ ninguem ”, “ esta ”, “ ate ”, “ tambem ” ).
Por outro lado, destaca -se o uso indevido do acento gráfico em palavras como “ melâncolia ” e “ intervênção ”, por exemplo.
Quanto ao domínio das convenções relativas à grafia das palavras, observaram -se desvios esporádicos como : hipercorreção pela escolha de “ e ” no lugar de “ i ”, por influência de hábitos da oralidade ( “ descriminação ” no lugar de “ discriminação ”, “ entervenção ” no lugar de “ intervenção ” ) ; desvios de grafia relacionados à variação diastrática, como por exemplo “ perca ”, “ subjulgadas ”, “ precoseito ”, “ soubre ”, “ vecendo ”, “ indesencia ”, “ apolojia ”, “ fulga ”, “ dereitos ” ;...
Observaram -se, também, muitos casos de inadequação no uso da maiúscula : uso indevido para destacar determinadas palavras-chave do texto, como “ Brasileiros ”, “ Homens ” ; ausência de maiúscula para grafar nomes próprios, como"marielle".
Expressões como"Direitos Humanos"foram grafadas diversificadamente, com maiúsculas e com minúsculas.
Destaque -se, também, a maioria dos participantes que grafam os textos inteiramente em caixa alta, o que impede a identificação de desvios relacionados ao uso de maiúsculas.
Vale observar, também, que, ao contrário do que se esperava, não apareceram abreviaturas próprias do “ internetês ”, relacionadas ao uso de redes sociais e e-mails, como por exemplo"vc".
Aspectos textuais : Para alguns dos graduandos essa competência se revelou como a mais problemática, tendo em vista os inúmeros problemas observados, desvios acumulados durante toda a formação escolar.
São eles : sequência justaposta de ideias sem encaixamentos sintáticos, provocando uma redução de estruturas subordinadas, ao lado do aumento na frequência de estruturas coordenadas e absolutas ; frases fragmentadas que comprometem a estrutura lógico-gramatical ; frases formadas apenas por oração subordinada, sem oração principal ; redução no uso de conectores para expressar relações lógicas essenciais à construção do texto como consequência da mudança de estruturação frasal ; emprego inadequado do pronome relativo ( com omissão da preposição ou a utilização do pronome relativo inadequado, como “ onde ” ), refletindo hábitos da oralidade ; ausência de recursos de referenciação, como substituição de termos por sinônimos, hiperônimos, nominalizações, expressões metafóricas.
Esses problemas revelam dificuldades em relação à estrutura formal do texto produzido, o que é preocupante ao se levar em conta que são graduandos em fase final de formação.
Em alguns textos, falta um mínimo de textualidade e de domínio do registro padrão da língua.
Na verdade, observaram -se relações linguísticas quase agramaticais, como as estabelecidas pela sequência de gerúndios sem o apoio de um ponto de partida para a organização das informações gramaticais e semânticas.
Quanto à utilização dos sinais de pontuação, observou -se uma ausência dos mesmos na maioria dos textos analisados.
Foi muito frequente a ocorrência de parágrafos sem marca interna de pontuação para separar os períodos.
São os seguintes os tipos de problemas mais encontrados :  vírgula : utilização de vírgula para separar o sujeito e o predicado ; uso de vírgula no lugar do ponto para separar ideias que constituem períodos distintos ; ausência de vírgula para separar oração adjetiva explicativa ;  ponto e vírgula : utilização do ponto e vírgula no lugar de vírgula ;  ponto final : ausência de ponto final para separar períodos.
Aspectos morfossintáticos e vocabulares : Os resultados são transparentes em relação aos aspectos mais problemáticos no desempenho dos participantes nestas duas competências.
Em relação aos aspectos morfossintáticos, seguem algumas observações sobre os desvios mais frequentes.
Quanto à regência, o desvio mais frequente foi a falta do sinal indicativo da crase – isso revela que o usuário não tem consciência de que, sob a forma do termo “ a ”, existe a presença de uma contração entre a preposição “ a ” ( exigida pela regência do termo anterior ) e o artigo definido “ a ”.
Outro problema relacionado à regência verbal e à nominal, encontrado frequentemente nas respostas, foi a ausência de preposição antes de pronome relativo, processo generalizado na modalidade oral da língua, em situações de registro informal.
Apesar da possibilidade de que essa alteração de regência se generalize no padrão escrito da Língua Portuguesa, como já está ocorrendo até em textos jornalísticos, o não emprego da preposição foi considerado como desvio neste processo de avaliação.
Outro desvio de regência significativo foi a utilização inadequada de uma preposição ou sua ausência após o verbo ou o nome ( substantivo ou adjetivo ).
A concordância verbal e a concordância nominal apresentaram alguns desvios muito frequentes.
Quanto à concordância de número, observou -se, como apontam as pesquisas nessa área, ausência de marca ( com sujeito anteposto ou posposto ) ou uso indevido ( uso inadequado da marca de plural comandado pelo núcleo plural da locução adjetiva, apesar de o substantivo que funciona como núcleo do sintagma nominal estar no singular ).
Uma ocorrência que se destacou foi a ausência de acento circunflexo na forma plural do presente do indicativo do verbo “ ter ”, que foi considerada como um desvio na concordância verbal e não na acentuação gráfica.
Quanto à concordância de gênero, vários casos foram observados, normalmente no âmbito de sintagmas nominais longos, em que o adjetivo está afastado do substantivo.
Quanto à questão da colocação pronominal, foram poucos os desvios observados.
Concluiu -se que, no registro escrito formal, a maioria dos participantes já incorporou regras como a não introdução da frase por um pronome oblíquo e a próclise na presença de um termo atrator.
Várias marcas de oralidade foram identificadas, embora não com alta frequência : uso do pronome relativo “ onde ” como relativo universal, falta de artigo definido antes de substantivo isolado ou antes de termos em enumerações.
Quanto aos aspectos vocabulares, alguns tipos de inadequação foram observados : expressões da oralidade ; seleção vocabular incompatível com o contexto, gerando falta de inteligibilidade ; falta de domínio de vocabulário mais abstrato e de maior complexidade, essencial ao desenvolvimento do texto de base dissertativa.
O principal aspecto observado foi a excessiva repetição de certas palavras, como o termo"pessoa", por exemplo, revelando limitação de repertório vocabular.
Em função do tema solicitado na Questão 1, diferentes grafias da expressão fake news não foram consideradas.
Para sintetizar, os problemas que mais se destacaram nos textos analisados foram : a ) truncamentos sintáticos, caracterizados por estruturas frasais incompreensíveis devido à ruptura da complexidade sintática própria da modalidade escrita ; b ) redução de estruturas subordinadas, compensada pelo aumento na frequência de estruturas coordenadas e absolutas, por um lado, ou pela elaboração de estruturas textuais prejudicadas pelo excesso de ideias sem a devida conexão subordinativa ; c ) redução no uso de conectores para expressar relações lógicas essenciais à construção do texto, substituídos pela exigência de inferência por parte do interlocutor para suprir a sua ausência ; d ) redução no uso do subjuntivo, ao lado da ampliação do uso do indicativo, combinado a estruturas frasais coordenadas ou absolutas ; d ) empobrecimento do processo de referenciação, com a repetição de palavras, sem processos mais sofisticados de substituição ; e ) simplificação da marcação da categoria tempo na morfologia verbal ; f ) desrespeito às regras de concordância verbal e nominal próprias do padrão formal ; g ) ausência do acento grave indicador da crase, revelando falta de conhecimento dos critérios de regência verbal e nominal ; h ) falta de domínio de vocabulário mais abstrato e de maior complexidade, essencial ao desenvolvimento do processo dissertativo ; i ) redução do emprego da acentuação gráfica.
Como resultado da correção, merece atenção a constatação de que a modalidade escrita tem apresentado a tendência a uma evidente simplificação, aproximando -se das características da modalidade oral da Língua Portuguesa.
No caso dos textos de base dissertativa ( expositiva ) e de base argumentativa, inscritos em um registro formal, a distância entre as duas modalidades é ainda maior, o que provoca desvios recorrentes em todos os quatro aspectos analisados : ortográficos, textuais, morfossintáticos e vocabulares.
Destaca -se, nessa avaliação, o comprometimento textual, por meio de estruturas fragmentadas e/ou truncadas, rompendo a complexidade sintática esperada no padrão formal. 6.3.2 Componente de Conhecimento Específico Na parte da prova relativa às questões discursivas do Componente de Conhecimento Específico ( Tabela 6.17 ), observa -se que a Média foi mais alta do que para as questões discursivas do Componente de Formação Geral.
Enquanto no Componente de Formação Geral, a Média para estudantes de Tecnologia em Design de Moda de todo o Brasil foi 36,2, na parte de Conhecimento Específico, a Média foi 42,8.
A maior Média desse componente foi obtida pelos estudantes da região Sul ( 43,9 ), e a menor, pelos da região Centro-Oeste ( 39,5 ).
Quanto à variabilidade das notas, o Desvio padrão de todo o Brasil foi 20,1.
O maior Desvio padrão foi encontrado na região Norte ( 20,8 ), e o menor, na região Sudeste ( 19,5 ).
A maior nota Máxima, 93,3, foi obtida nas regiões Nordeste e Sudeste.
Nas demais regiões, as notas Máximas foram : 83,3 na região Norte, 86,7 na Sul e 90,0 na região Centro-Oeste.
A nota Mínima ( 0,0 ) foi obtida por, pelo menos, um aluno em todas as regiões do Brasil.
A Mediana do Brasil, como um todo, foi 45,0, obtida também nas regiões Norte e Nordeste.
Na região Sudeste, a Mediana foi 41,7.
A maior Mediana foi obtida na região Sul ( 46,7 ), e a menor, na região Centro-Oeste ( 40,0 ), ou seja, na região Centro-Oeste pelo menos metade dos alunos tirou nota 40,0 nas questões discursivas do Componente de Conhecimento Específico.
No Gráfico 6.10, representa -se a distribuição das notas nas questões discursivas no Componente de Conhecimento Específico.
A moda dessa distribuição ocorre no intervalo ( 40 ; 50 ], com quase 20 \% do total de participantes, seguido de perto pelos intervalos ( 50 ; 60 ] e ( 30 ; 40 ].
A análise de cada uma destas questões será feita a seguir. 6.3.2.1 Análise da Questão Discursiva 3 do Componente de Conhecimento Específico Na Questão 3, cujos resultados aferidos se encontram descritos na Tabela 6.18, a Média dos estudantes de todo o Brasil foi 52,5, a questão de melhor desempenho dentre as três discursivas de Conhecimentos Específicos.
A menor Média nessa questão foi obtida pelos alunos da região Centro-Oeste ( 43,6 ), enquanto a maior Média foi obtida na região Nordeste ( 56,8 ).
Quanto à variabilidade das notas, o Desvio padrão de todo o Brasil foi 28,3.
O maior Desvio padrão foi obtido na região Norte ( 34,2 ), enquanto o menor foi obtido na região Sul ( 26,8 ).
A nota Máxima, 100,0 pontos, foi alcançada por, pelo menos, um aluno de quase todas as regiões, com exceção da região Centro-Oeste que teve nota Máxima 90,0.
A Mediana do Brasil, como um todo, foi 60,0, a mesma para três regiões.
Nas regiões Sul e Centro-Oeste, a Mediana foi 50,0.
A nota Mínima ( 0,0 ) também foi a mesma em todas as regiões do Brasil.
No Gráfico 6.11, mostra -se a distribuição das notas na Questão discursiva 3, do Componente de Conhecimento Específico da área de Tecnologia em Design de Moda.
A moda dessa distribuição é o intervalo ( 70 ; 80 ], seguido de perto pelos intervalos ( 60 ; 70 ] e ( 50 ; 60 ]. 6.3.2.2 - Comentários sobre as respostas à Questão Discursiva 3 A Questão discursiva 3 da área de Tecnologia em Design de Moda versava sobre a adequação do vestuário e seus acessórios complementares às necessidades dos diferentes No Gráfico 6.11, mostra -se a distribuição das notas na Questão discursiva 3, do Componente de Conhecimento Específico da área de Tecnologia em Design de Moda.
A moda dessa distribuição é o intervalo ( 70 ; 80 ], seguido de perto pelos intervalos ( 60 ; 70 ] e ( 50 ; 60 ]. 6.3.2.2 - Comentários sobre as respostas à Questão Discursiva 3 A Questão discursiva 3 da área de Tecnologia em Design de Moda versava sobre a adequação do vestuário e seus acessórios complementares às necessidades dos diferentes tipos físicos dos usuários, e que esse mercado, o da Moda Inclusiva, vem \textbf{crescendo} gradualmente e sendo valorizado ainda mais nos dias hoje.
Nesse contexto, durante o desenvolvimento de produtos de moda, os conhecimentos relacionados à ergonomia e a metodologia projetual são primordiais para o atendimento das necessidades desse mercado.
A questão estava clara e objetiva e destacou a moda inclusiva como tema a ser abordado, principalmente, tendo em vista a ênfase sobre a ergonomia.
Nesse sentido, a temática da questão compreendeu as Diretrizes Curriculares para a área de Tecnologia em Design de Moda, quando evidenciou o perfil"crítico, reflexivo e ético quanto aos aspectos sociais, ambientais, políticos, históricos, culturais, educacionais e inclusivos que impactam a área da moda".
A abordagem da ergonomia como área de conhecimento específica ao design de moda sugeriu a relação entre corpo e vestuário e direcionou a atenção à competência de"criar e desenvolver produtos e/ou coleções de moda viáveis, de acordo com a metodologia projetual de design e de moda".
Apesar dessa direção, percebeu -se que ela também evidenciou uma perspectiva sobre um perfil do estudante comprometido com os fatores ergonômicos e sua aplicação, o que caracteriza as etapas de desenvolvimento de produtos de moda e reforça a associação da questão à competência de aplicação da metodologia projetual da área.
A ergonomia aplicada à moda inclusiva demandou do aluno uma capacidade de interpretação das necessidades do usuário, nesse caso em relação ao produto do vestuário, fazendo com que haja uma adequação na definição e na indicação de soluções a partir de problemas específicos.
Dessa forma, o nível de profundidade que a questão exigiu foi de conhecimento médio sobre o tema, pois a questão estava fundamentada em conteúdos que não correspondem àqueles mais tradicionais e mais difundidos no campo da moda, como é o caso do desenvolvimento de coleção.
Um ponto importante a ser observado é que o enunciado da questão aborda a Moda Inclusiva de uma forma bastante ampla, citando diferentes tipos de deficiência, como mental, auditiva e física.
Essa estratégia discursiva da questão desempenhou um papel de seleção dos conteúdos e de verificação da capacidade interpretativa do aluno, uma vez que a própria definição de deficiência foi posta em questionamento, na medida em que a pergunta solicitou sua identificação.
Esses fatores provocaram a avaliação das condições de entendimento do comando e também sobre o conhecimento a ser empregado para a resolução da questão.
As etapas previstas pelo comando, no que concerne à citação da deficiência, aos problemas e as adaptações do vestuário, foram contempladas nas repostas dos alunos a partir do emprego de uma linguagem objetiva e adequada aos termos da ergonomia e ao vocabulário do vestuário e da área de Tecnologia em Design de Moda. partir do emprego de uma linguagem objetiva e adequada aos termos da ergonomia e ao vocabulário do vestuário e da área de Tecnologia em Design de Moda.
A maioria das respostas realizou abordagens corretas dos problemas de vestuários e das adaptações possíveis para solucionar os problemas tendo em vista a especificidade da deficiência.
No entanto, apesar das abordagens corretas, elas não corresponderam ao nível de detalhamento previsto, principalmente, em relação à descrição dos problemas do vestuário.
Nesse sentido, as repostas demonstraram atender parcialmente aos critérios de avaliação, já que a generalização da resposta, frequentemente, se associava ao problema do usuário, não alcançando a particularidade do vestuário.
Uma possível causa para esse fator pode estar na desatenção ao comando da questão e, portanto, à especificidade da categoria da “ vestibilidade ” associada à deficiência apontada pelo aluno.
Ainda sobre esse aspecto, foi possível observar que as respostas que desenvolveram parcialmente a descrição dos problemas indicaram a “ vestibilidade ” como problema, principalmente, através do ato de vestir e despir.
Nesse sentido, percebe -se uma associação desse tipo de resposta à indicação do texto de abertura do enunciado da questão, já que nele havia a afirmação de que"os principais objetivos da Moda Inclusiva são simplificar o ato de vestir e despir".
A maior incidência de respostas equivocadas foi na citação das deficiências.
As respostas dos estudantes apresentaram indicações designando doenças e distúrbios como deficiências.
Nesse sentido, o que parece ter ocorrido é que as afirmações se associaram à amplidão do sentido de Moda Inclusiva que o texto do enunciado da questão declarava, estabelecendo, assim, pouca reflexão sobre a especificidade do sentido de deficiência.
Da mesma forma, essa constatação permite pensar que a própria definição de deficiência tenha gerado dúvida, já que ela possui uma diversidade de exemplos e que, compreendem, inclusive, conceituações e debates da área da saúde.
Vale a pena ressaltar que também foram encontradas respostas que indicaram deficiências que não representavam casos de problemas de “ vestibilidade ”, o que, embora sugerisse o atendimento correto à demanda de citação da deficiência, se tornava uma incoerência em relação ao objetivo da questão pela perspectiva da área de conhecimento do design de moda e, em especial, da ergonomia aplicada ao vestuário.
As deficiências mais citadas foram a física e a visual, através da indicação dos cadeirantes e cegos.
No caso dos cadeirantes, os problemas se concentraram no ato de vestir as peças do vestuário, em específico, as partes de baixo.
Essa sinalização demonstrou uma abordagem parcial do problema, já que poucas vezes foram citadas as causas no vestuário para os problemas de vestir e/ou despir.
Em relação às soluções para os problemas dos cadeirantes, a resposta mais aplicada foi a colocação de sistemas de fechamentos, como abotoamentos ou velcros, nas laterais das calças para facilitar o ato de se vestir sozinho.
Também foram comentados os problemas de posicionamento de aviamentos e o contato intermitente com a pele, devido à imobilidade do corpo.
Nos casos em que a cegueira foi apontada como deficiência, houve a menor variação nas respostas em relação ao padrão de resposta.
Comumente, o problema de identificação das roupas ( cores, formas e estampas ) foi apresentado nos textos e o emprego de etiquetas em braile, o uso de texturas e relevos foram citados como a solução adequada.
A diversidade de respostas também contemplou o nanismo.
Em geral, quando mencionavam a deficiência, as respostas abordavam a questão da dificuldade do usuário de encontrar roupas para seus corpos, exemplificando através da compra em lojas infantis.
Para solucionar esse problema, os estudantes indicaram o ajuste da roupa ao corpo, o que configurava como uma resposta genérica, sem detalhar as relações de proporção e medidas do corpo com a roupa.
Dentre as respostas equivocadas, uma parcela pareceu seguir a contextualização da questão e apontou públicos-alvo de produtos de moda inclusiva, mas que não são enquadrados como deficiências.
A saber : plus size ( obesidade ), gestantes, depressão, entre outros.
Nesse mesmo sentido, algumas repostas citaram a deficiência auditiva e apontaram problemas de “ vestibilidade ” e soluções, o que se configurou como um grande erro, tendo em vista que esse público não se configurava pelo objetivo daquele identificado na questão.
A maioria das respostas obteve um desenvolvimento parcial dos objetivos da questão, apresentando, em suas abordagens, uma variação na forma como os problemas se relacionavam com as deficiências ou no modo como as adaptações sugeridas se adequavam aos problemas citados.
As notas baixas restringiram a resposta às citações das deficiências, ou à apresentação da deficiência com indicação do problema e/ou da adaptação incipiente ou, até mesmo, inadequado ao contexto.
Dentre as respostas com notas altas, percebeu -se pertinência nas relações apontadas entre as deficiências, os problemas e as adaptações, apresentando, algumas vezes, leves inadequações na descrição do problema ou na eficiência da adaptação.
O enfoque da Questão 3 foi sobre um tema específico do projeto em design de moda e a ergonomia aplicada aos casos das deficiências, contudo essa temática não compreende os enfoques comumente abordados no desenvolvimento de projetos de moda, uma vez que a coleção é um formato mais tradicional do mercado, e, portanto, mais enfatizado na formação dos estudantes.
Nesse sentido, o nível médio das respostas pareceu representar uma visão panorâmica do conhecimento dos estudantes, principalmente, devido às generalizações das respostas, demonstrando a falta de uma abordagem mais empírica no tratamento das informações.
As generalizações, nesse contexto, poderiam indicar pouca profundidade em relação ao desenvolvimento de projetos com esse enfoque no âmbito do ensino.
Tendo em vista a forma como os problemas solicitados pela questão foram descritos, percebeu -se uma limitação nas capacidades analíticas e de pesquisa dos estudantes, em especial, aquelas que são exigidas em contato com o público.
Tal fator sugere inconsistências nas metodologias de projeto adotadas, o que também poderia ser reflexo da pouca bibliografia e de pesquisas sobre o tema que favoreçam a difusão do conhecimento.
Os problemas terminológicos e de conceituação das deficiências que apareceram nas respostas põem em pauta a falta de uma perspectiva mais definida sobre o trabalho em design e, principalmente, sobre sua condição interdisciplinar, visto que a questão abordava conceitos e definições da área da saúde.
As soluções propostas pelos estudantes também sinalizaram uma falta de reconhecimento das estruturas materiais dos produtos de vestuário, evidenciando uma lacuna no conhecimento prático sobre o objeto da área de formação.
As funções das partes do vestuário parecem ser desconhecidas, o que põe em pauta uma aprendizagem pouco concreta e que não favorece a experiência material do trabalho de designer de moda.
Assim, a correção da Questão 3 permitiu considerar que o perfil reflexivo e crítico do estudante provocado pela questão precisa ser apoiado também em um comprometimento e uma responsabilidade sobre os fatores de desenvolvimento dos produtos na sociedade em que o design de moda se insere.
Já que o tema se refere à moda inclusiva e aos aspectos da ergonomia, é necessário pôr em destaque os aspectos materiais pelos quais os usuários se relacionam com os objetivos de modo a direcionar a atuação do designer para a identificação das necessidades e para propostas de soluções efetivas e inovadores frente ao cenário industrial e padronizado do campo da moda. 6.3.2.3 - Análise da Questão Discursiva 4 do Componente de Conhecimento Específico Na Tabela 6.19, constam as informações relativas à Questão 4 do conjunto de questões do Componente de Conhecimento Específico.
O desempenho dos estudantes nessa questão foi inferior ao das questões 3 e 5.
A Média geral do Brasil foi 37,7, sendo a maior Média registrada na região Sul ( 40,2 ), e a menor, na região Nordeste ( 35,6 ).
A nota Máxima ( 100,0 ) foi atingida por, pelo menos, um concluinte em quatro regiões, exceto na Norte, com nota Máxima 60,0.
A Mediana em todo o Brasil foi 40,0, o mesmo valor foi obtido em quase todas as regiões, exceto na Centro-Oeste, com Mediana 37,5.
A nota Mínima de todas as regiões foi zero.
No Gráfico 6.12, está representada a distribuição de notas da Questão discursiva 4, do Componente de Conhecimento Específico.
Essa distribuição tem moda na classe de estudantes do intervalo ( 30 ; 40 ], seguida de perto pelos intervalos ( 40 ; 50 ] e ( 50 ; 60 ].
O intervalo ( 10 ; 20 ] caracteriza -se como um máximo local, com cerca de 12 \% do total. 6.3.2.4 Comentários sobre as respostas à Questão Discursiva 4 A Questão 4 continha 2 textos introdutórios, um que versava sobre a importância da pesquisa, da disciplina e das técnicas necessárias para auxiliar o profissional de design de moda no processo de criação ; e outro que trazia uma breve descrição sobre os moodboards ou storyboards e de sua importância na tarefa de reunir informações e na apresentação de ideias.
O comando da questão solicitava a apresentação de duas contribuições dos painéis semânticos ( moodboards ou storyboards ) para o processo de criação de produtos de moda, além da indicação de dois tipos de pesquisa que poderiam ser apresentadas em modelo de painel semântico e para cada tipo de pesquisa citada a apresentação de dois elementos que poderiam compor o painel.
De acordo com as Diretrizes Curriculares para a área de Tecnologia em Design de Moda, a Questão 4 se apoia nas competências de"Criar e desenvolver produtos e/ou coleções de moda viáveis, de acordo com a metodologia projetual de design e de moda".
Pelo enunciado da questão, observou -se que o enfoque direcionou ao tema das metodologias projetuais, ressaltou a relação entre as pesquisas que os designers de moda realizam e a função dos painéis semânticos como recursos para as considerações destas pesquisas.
O perfil de referência da questão sugeriu um profissional"criativo, inovador e empreendedor para atuar nos diversos segmentos da moda"e destacou a pesquisa como componente específico dos conteúdos de referência da área de atuação.
Os conhecimentos exigidos na questão representaram um fundamento básico desenvolvido nos cursos da área, tendo em vista que exigiu do aluno um nível de profundidade elementar.
A pesquisa e o painel são procedimentos constantes em todas as atividades de projeto que os estudantes venham a realizar durante o curso, o que o aproxima, inclusive, das atividades estipuladas nas suas atribuições profissionais no mercado de trabalho.
A Questão 4 pode ser considerada uma questão fácil para os estudantes.
O conteúdo selecionado e a forma de abordagem se enquadraram em noções muito gerais da formação em Tecnologia em Design de Moda, inclusive, que são empregadas para o desenvolvimento de níveis avançados da atividade profissional.
Por se tratar de uma questão com dois itens, sendo que para cada um deles havia objetivos quantitativos de informação distintos a ser respondidos, as respostas compreenderam uma diversidade de formas de composição dos elementos que deviam ser contemplados no texto, apresentando, assim, formulações incompletas.
Embora a forma de abordagem das respostas fosse incompleta, a maioria apresentou formulações corretas.
Em relação ao item ‘ a ’, as contribuições dos painéis atingiram a diversidade de respostas esperada, porém, algumas vezes, o número de contribuições era inferior àquele solicitado no enunciado da questão.
Vale a pena ressaltar que, nesse item, a identificação da ação de parafrasear o enunciado da questão foi dificultada devido à proximidade do texto do enunciado da questão e da resposta considerada como correta, o que configurou uma facilidade para as respostas dos estudantes.
No item ‘ b ’ da questão, foi encontrada a maior parte de respostas equivocadas.
Na citação dos tipos de pesquisa e dos elementos que poderiam compor o painel, surgiram respostas que mencionavam pesquisas que não estavam contempladas nas metodologias projetuais do design de moda, ora empregando uma terminologia proveniente das metodologias científicas, ora designando as pesquisas pelas suas fontes.
Pelo caráter circular das informações na questão, as respostas podem ter -se direcionado a esses conteúdos por uma dúvida sobre a validade de repetição do conteúdo de resposta para os dois itens.
O comando da questão permitia o emprego de informações muito similares para a composição das respostas dos itens, o que também era reforçado pelos textos de abertura do enunciado da questão.
Para o aluno, essa redundância pode ter representado a necessidade de escolha de um caminho diferente, levando -o a pensar que talvez a resposta não poderia ser tão semelhante.
É necessário destacar que algumas vezes os títulos empregados para designar as pesquisas eram frutos de adaptações sobre os títulos utilizados pela referência bibliográfica, como é o caso da pesquisa de inspiração, cujo título não apresentava a designação de"tema da coleção", mas em sua explicação, os alunos mencionavam o tema.
Outro exemplo é a variação de indicação da pesquisa de comportamento, que foi encontrada como pesquisa de público-alvo e pesquisa de comportamento do consumidor.
A maioria das respostas no item ‘ b ’ não apresentou os elementos dos painéis conforme solicitado no enunciado da questão ; apenas os nomearam, sem distinguir sua relação com a pesquisa.
Os tipos de respostas mais encontrados seguiram uma abordagem de acordo com a divisão proposta pela questão, isto é, uma separação entre o Item ‘ a ’ e o item ‘ b ’.
Poucos estudantes construíram respostas que apresentaram as informações combinadas dos itens.
Quando o fizeram, acentuaram a circularidade dos dados abordados.
Em relação às contribuições, as mais mencionadas foram a de"guiar o processo de criação"e a de reunir informações de tendências.
Nesse sentido, as respostas destacaram, principalmente, os elementos criativos, como cores, tecidos, dentre outros que aludiam diretamente ao Texto 2 do enunciado da questão.
Muitas respostas citaram como contribuição dos painéis o procedimento de criação de extração das cores.
É interessante observar que essa designação referência uma ação muito específica do trabalho de criação, estabelecida principalmente nas metodologias de ensino do design de moda e que emprega o painel como inferior àquele solicitado no enunciado da questão.
Vale a pena ressaltar que, nesse item, a identificação da ação de parafrasear o enunciado da questão foi dificultada devido à proximidade do texto do enunciado da questão e da resposta considerada como correta, o que configurou uma facilidade para as respostas dos estudantes.
No item ‘ b ’ da questão, foi encontrada a maior parte de respostas equivocadas.
Na citação dos tipos de pesquisa e dos elementos que poderiam compor o painel, surgiram respostas que mencionavam pesquisas que não estavam contempladas nas metodologias projetuais do design de moda, ora empregando uma terminologia proveniente das metodologias científicas, ora designando as pesquisas pelas suas fontes.
Pelo caráter circular das informações na questão, as respostas podem ter -se direcionado a esses conteúdos por uma dúvida sobre a validade de repetição do conteúdo de resposta para os dois itens.
O comando da questão permitia o emprego de informações muito similares para a composição das respostas dos itens, o que também era reforçado pelos textos de abertura do enunciado da questão.
Para o aluno, essa redundância pode ter representado a necessidade de escolha de um caminho diferente, levando -o a pensar que talvez a resposta não poderia ser tão semelhante.
É necessário destacar que algumas vezes os títulos empregados para designar as pesquisas eram frutos de adaptações sobre os títulos utilizados pela referência bibliográfica, como é o caso da pesquisa de inspiração, cujo título não apresentava a designação de"tema da coleção", mas em sua explicação, os alunos mencionavam o tema.
Outro exemplo é a variação de indicação da pesquisa de comportamento, que foi encontrada como pesquisa de público-alvo e pesquisa de comportamento do consumidor.
A maioria das respostas no item ‘ b ’ não apresentou os elementos dos painéis conforme solicitado no enunciado da questão ; apenas os nomearam, sem distinguir sua relação com a pesquisa.
Os tipos de respostas mais encontrados seguiram uma abordagem de acordo com a divisão proposta pela questão, isto é, uma separação entre o Item ‘ a ’ e o item ‘ b ’.
Poucos estudantes construíram respostas que apresentaram as informações combinadas dos itens.
Quando o fizeram, acentuaram a circularidade dos dados abordados.
Em relação às contribuições, as mais mencionadas foram a de"guiar o processo de criação"e a de reunir informações de tendências.
Nesse sentido, as respostas destacaram, principalmente, os elementos criativos, como cores, tecidos, dentre outros que aludiam diretamente ao Texto 2 do enunciado da questão.
Muitas respostas citaram como contribuição dos painéis o procedimento de criação de extração das cores.
É interessante observar que essa designação referência uma ação muito específica do trabalho de criação, estabelecida principalmente nas metodologias de ensino do design de moda e que emprega o painel como Os diferentes tipos de informações a serem reunidas como contribuição dos painéis, também foram citados, mas com menos frequência em relação às tendências.
Em uma ordem quantitativa, as informações de público eram seguidas daquelas sobre a marca.
A função de ilustração foi mencionada em algumas respostas, mas ela se apresentava com frequência naquelas mais elaboradas e que conseguiam, inclusive, fazer uma distinção entre duas contribuições diferentes.
Na maioria das respostas, as informações utilizadas para responder ao item ‘ a ’, na verdade, representava uma única contribuição, como foi o caso do sentido de “ reunir informações ”, para o qual somente era listado o tipo de informação e sugerido pela indicação no enunciado da questão.
O principal equívoco cometido para a resposta do item ‘ a ’ foi a citação dos elementos de criação como cores, tecidos, aviamentos, entre outros, de uma forma muito objetiva e pontual.
Em relação ao item ‘ b ’, as pesquisas mais mencionadas foram a de tendências, a tema de coleção e a de comportamento.
Quando apresentados, os elementos da pesquisa de tendência e de tema da coleção eram tecidos, cores, formas e texturas.
No caso da pesquisa de comportamento, eram citados lugares que o público frequenta, seus estilos e hábitos de consumo.
Muitas respostas não chegaram a designar os elementos dos painéis.
Outras, empregaram diferentes termos para aludir às pesquisas.
Os termos"pesquisa/painel de ambiência"foram encontrados com frequência para designar a pesquisa e o painel de tema da coleção.
A consideração como resposta foi possível nos casos em que o estudante realizou a descrição da pesquisa e, assim, empregou termos como tema e coleção em sua explicação.
A pesquisa de comportamento do consumidor também apresentou diferentes designações, como público-alvo, comportamento do consumidor e estilo de vida.
Algumas respostas utilizaram a expressão"Pesquisa de moda"para se referir às tendências e outras fizeram separação entre pesquisa de público-alvo e de estilo de vida.
Em relação aos equívocos, percebeu -se que eles se concentraram na identificação das pesquisas.
Algumas respostas empregaram nomes de pesquisas referentes à metodologia da pesquisa, citando expressões como"pesquisa visual","imagética","bibliográficas".
Em alguns casos, os elementos citados eram associados a fontes de pesquisa : livros, revistas, internet, entre outros.
Frequentemente foi empregada, equivocadamente, na designação das pesquisas a distinção entre macrotendências e de microtendências.
Isso não estava de acordo com a referência bibliográfica, havendo grande variação de definição sobre estas expressões.
Alguns estudantes em suas respostas cometeram equívocos ao exemplificarem os elementos dos painéis, citando possibilidades de temas e suas imagens, como por exemplo, um painel cujo tema seria"natureza", poderia conter"flores".
A correção da questão permitiu visualizar algumas especificidades sobre o conteúdo abordado.
Em primeiro lugar, é possível citar a forma como a maioria das respostas referentes às contribuições dos painéis demonstrou uma limitação na compreensão da função deles na metodologia projetual, principalmente, pelo fato de sinalizarem restrições ao seu papel informacional e criativo.
A pouca indicação da função ilustrativa do painel sugere uma falta de clareza ao fato de ele ser uma ferramenta de transformação de informações textuais para a dimensão visual.
Sua contribuição ao reunir informações deixa evidente que o ato de reunir já é um processo de tradução dos dados em diferentes linguagens.
Assim, a pouca evidência dessa característica parece sinalizar uma incompreensão da constituição visual que a metodologia projetual do design de moda pressupõe e que é inerente ao próprio trabalho da moda, já que seu suporte é a aparência.
Em segundo lugar, a variação apresentada na designação das pesquisas que o designer de moda realiza indica que há a sobreposição de termos e, portanto, de conteúdos, principalmente, em relação ao papel da pesquisa científica no campo do design de moda.
O emprego de termos referentes a tipos de pesquisas e fontes para as respostas do item ‘ b ’ demonstraram que as semelhanças dos termos criaram uma confusão entre metodologia de pesquisa e metodologia projetual em design de moda.
A distinção do ato de pesquisar na prática profissional do designer em relação à pesquisa realizada para fins de um trabalho acadêmico e científico deve ser realizada para que a própria metodologia projetual seja entendida em sua diversidade e associada ao âmbito da prática de desenvolvimento de produto, cujos painéis são ferramentas.
Ainda sobre os termos empregados pelos estudantes, há outro aspecto que as respostas parecem revelar sobre a formação na área : os regionalismos.
Durante a correção foi possível identificar blocos de respostas que empregam diferentes expressões para designar os tipos de pesquisa, o que sugere uma variação setorizada sobre termos que são recorrentes e bastante alinhados na bibliografia geral do campo do design, do qual a moda advém.
Assim, esses regionalismos podem evidenciar pouco alinhamento sobre as referências bibliografias da área e, em especial, do conteúdo abordado na questão.
Sabe -se que a metodologia projetual do design de moda não possui vasta bibliografia no Brasil, porém, o título empregado pela questão corresponde a um dos mais tradicionais e difundidos sobre o tema.
Nesse sentido, a ausência de alinhamento dos termos parece apontar para formações desassociadas das próprias diretrizes da área.
Pelo desempenho dos estudantes, verificou -se que o perfil crítico e inovador sugerido pela questão não foi atingido pelas respostas.
De um modo geral, o desempenho refletiu uma condição bastante restrita de compreensão do projeto em design, principalmente, pela maneira como as contribuições dos painéis e das etapas de pesquisa foram identificadas.
Pela alusão ao trabalho criativo que a questão sugere, as respostas parecem confirmar uma formação voltada aos \textbf{parâmetros} industriais de desenvolvimento de coleção e que correspondem à \textbf{reprodução} de estratégias estéticas como aquelas associadas aos temas das coleções e às tendências.
Portanto, o desempenho sinaliza a possibilidade de enfoques mais amplos e que permitam a compreensão da metodologia projetual e a pesquisa associada às especificidades da área de design de moda. 6.3.2.5 Análise da Questão Discursiva 5 do Componente de Conhecimento Específico Na Tabela 6.20, constam as informações relativas à Questão 5 do conjunto do Componente de Conhecimento Específico.
O desempenho dos estudantes de todo o Brasil nesta questão foi inferior ao da Questão 3 e superior ao da Questão 4 dentre as três questões discursivas desse componente.
A nota Média dos estudantes de todo o Brasil foi 38,2.
A maior Média foi registrada na região Sul ( 40,6 ), enquanto a menor Média foi registrada na região Norte ( 30,0 ).
Quanto à variabilidade das notas, o Desvio padrão dos alunos do Brasil, como um todo, foi 29,5.
Enquanto o maior desvio foi encontrado na região Sul ( 30,5 ), o menor foi encontrado na região Norte ( 23,9 ).
A Mediana para o Brasil e na maioria das regiões foi 25,0, a mesma das regiões Norte, Nordeste e Sudeste.
Nas demais regiões, a Mediana foi 50,0.
Para o conjunto de alunos de Tecnologia em Design de Moda do Brasil, a nota Máxima foi 100,0 e a nota Mínima, zero, em todas as regiões.
No Gráfico 6.13, é representada a distribuição das notas da Questão discursiva 5 do Componente de Conhecimento Específico.
Essa distribuição tem o intervalo modal ( 20 ; 30 ], seguido de perto pelo ( 40 ; 50 ], concentrando mais da metade dos participantes. 6.3.2.6 Comentários sobre as respostas à Questão Discursiva 5 A Questão 5 introduziu o tema da sustentabilidade através da descrição do caso real, de um estilista que se propôs a ajudar um grupo de artesãs que desejava utilizar os restos de material de um polo têxtil.
Esse estilista já criava coleções de acessórios e peças de decoração aproveitando os resíduos de materiais e dessa forma, produzia novos produtos.
O texto introdutório apresentava linguagem objetiva e narrativa contextualizando o estudante no tema da sustentabilidade, permitindo a comunicação do comando da questão sem gerar dubiedade.
O tema da questão se insere nos conteúdos de"sustentabilidade aplicada ao desenvolvimento de produtos e serviços de moda", conforme citam as Diretrizes Curriculares da Área de Tecnologia em Design de Moda.
Nesse sentido, a sustentabilidade se direciona à competência"empreender e gerenciar negócios de moda", estimulando um perfil discente"crítico, reflexivo e ético quanto aos aspectos sociais, ambientais, políticos, históricos, culturais, educacionais e inclusivos que impactam a área da moda".
Também foi possível verificar que a sustentabilidade incentiva um profissional responsável e comprometido em sua atuação profissional, principalmente, no âmbito da gestão e no desenvolvimento de produtos, expandindo as competências às ações de avaliação dos processos produtivos.
Analisando o nível da questão, pode -se classificar a questão como sendo de nível médio.
A Questão 5 exigiu do estudante uma capacidade de análise, tendo em vista que o enunciado demandava a descrição das especificidades e das dificuldades da indústria no âmbito da sustentabilidade, e uma capacidade de síntese, já que o estudante deveria propor soluções que demonstrassem reflexão sobre o tema associado ao design de moda.
Nesse sentido, percebeu -se uma abrangência para a composição da resposta, o que favoreceu as possibilidades de interpretação.
Em específico, foi possível estabelecer a relação entre a dificuldade e a solução solicitada no comando da questão, de modo objetivo, no processo de avaliação das repostas.
Entretanto, o caso narrado no enunciado da questão oferecia indícios para a resposta, por delimitar um tipo de dificuldade e uma solução, restringindo a diversidade de intepretações sobre o tema e as possibilidades de resposta da questão.
A dificuldade “ separação e descarte das sobras de tecido e linha ” e suas respectivas soluções : “ encaminhamento para coleta seletiva ; utilização na criação de novos produtos ( upcycling ) ; e utilização em outros segmentos da indústria : construção civil, automobilística, naval, entre outras ” ; se assemelhavam ao caso na citação de abertura, o que induziu as possibilidades de respostas dos estudantes, comprometendo a função \textbf{avaliativa} da questão.
Os textos dos estudantes apresentaram variadas composições de respostas, contemplando diferentes dificuldades e soluções.
A maioria dos estudantes utilizou de objetividade para responder à questão, criando associações diretas entre as dificuldades e as soluções, entretanto, sem o nível de especificidade esperado em relação aos produtos, elementos ou até mesmo malefícios indicados na descrição dos problemas.
As respostas privilegiaram problemas descritos superficialmente, indicando soluções genéricas que demonstravam pouca abrangência no âmbito da sustentabilidade.
A pouca profundidade na apresentação das dificuldades e também das soluções pode ter sido causada pela falta de familiaridade com a realidade industrial, o que é proporcional ao contexto de formação dos discentes e sua inserção no mercado profissional.
Apesar dessa distinção no detalhamento das respostas, foram identificados alguns problemas contextualizados no debate da sustentabilidade na indústria têxtil e de confecção. principalmente, em relação aos nomes de substâncias, processos ou ferramentas de trabalho, como upcycling, zero waste, biodegradável, agrotóxicos, logística, economia circular, entre outros.
No entanto, com o uso da linguagem mais genérica, as respostas conseguiram comunicar o sentido amplo daquilo que estavam designando.
Na diversidade de respostas encontradas, as equivocadas foram resultado de incoerências entre as dificuldades e as soluções, de modo que não se estabelecia um atendimento ao critério de pertinência da solução à dificuldade.
As respostas incompletas compreenderam casos nos quais os estudantes não atenderam ao objetivo quantitativo do comando : em vez de citarem 2 dificuldades e suas soluções, só citaram uma.
Em alguns casos, foram percebidas citações de dificuldades semelhantes, quase idênticas, que o aluno apresentava como se fossem distintas.
Em outras respostas, alunos apresentaram mais de uma solução para a mesma dificuldade.
Houve casos em que os alunos somente apresentaram dificuldades, sem propor as soluções.
Cabe destacar que o exemplo do caso citado no enunciado da questão configurou a construção de algumas respostas dos estudantes.
O problema de sobra de tecidos e do descarte de resíduos foi descrito com frequência como dificuldade e o encaminhamento do material a ONGs para reaproveitamento, bem como sua utilização para artesanato, como solução.
O direcionamento de resíduos a ONGs e para a produção de artesanato sugere, implicitamente, a criação de novos produtos.
Assim, esse tipo de resposta alude diretamente ao exemplo utilizado pelo enunciado, uma vez que as artesãs são citadas como uma alternativa de parceria com designers para a criação de novos produtos.
Essa constatação demonstra a efetivação de um direcionamento causada pelo enunciado da questão e gera uma restrição na dimensão \textbf{avaliativa} da questão.
A abordagem mais frequente para responder à pergunta foi aquela que identificou as dificuldades e as suas soluções a partir de dois casos : o descarte de resíduos e o beneficiamento de tecidos.
A resposta mais citada contemplou a categoria do descarte, mencionando a sobra de tecidos de coleções anteriores, bem como os resíduos provenientes das etapas de corte e confecção.
Essa ocorrência fortalece a indução realizada pelo enunciado da questão e demonstra uma restrição referente à contextualização do problema.
A segunda dificuldade tratou sobre as consequências do beneficiamento têxtil, citando, em especial a indústria de jeans, os efeitos no meio-ambiente de suas lavagens e a poluição da água.
Vale destacar que o jeans estava citado no enunciado da questão, já que o exemplo mencionado era sobre os resíduos deste tipo de material.
Acerca das soluções apontadas a partir destes problemas, o direcionamento de corte, foram as mais citadas.
As respostas que demonstraram maior desenvolvimento empregaram termos como reaproveitamento, upcycling e o zero waste, evidenciando uma profundidade nas propostas.
Porém, as respostas mais superficiais, muitas vezes empregaram o recurso da exemplificação, sugerindo a confecção de um novo produto a partir do material descartado.
Nesse sentido, houve uma recorrência para a citação dos retalhos de tecidos de malha para a confecção de tapetes.
Esses casos se associaram novamente àquele do enunciado da questão, já que a imagem ilustrava um produto gerado pela trama dos retalhos de jeans.
Para as respostas que abordavam as dificuldades geradas pelos beneficiamentos de tecidos, foram propostas soluções que envolviam o tratamento da água, incluindo novos processos de lavagem no caso do jeans, além do emprego de novos materiais, como aqueles de origem orgânica.
Em relação a essas respostas, o desempenho dos estudantes chegou mais próximo o esperado, evidenciando uma especificidade no tema da questão.
Uma das propostas mais mencionadas pelos estudantes para solucionar diferentes dificuldades foi a conscientização através de campanhas de educação.
Muitas respostam citaram campanhas para os fabricantes, para os designers, para os consumidores, identificando esse recurso para diferentes categoriais de dificuldades, como a origem dos insumos, o descarte, beneficiamento, processo produtivo e pós-consumo.
Alguns estudantes descreveram dificuldades que não estavam contextualizadas no debate da sustentabilidade, como era o caso da falta de mão de obra, o excesso de trabalho e a aceleração da produção.
Embora esses problemas sejam específicos da indústria têxtil e de confecção, da forma genérica como foram apresentados, eles não eram contemplados pela perspectiva mais específica, salvo quando a resposta do estudante se encaminhava para uma associação destas dificuldades ao universo do mundo da moda, como por exemplo, a citação do fast fashion.
Outra parte das respostas não alcançou a relação entre as dificuldades e as soluções, revelando, muitas vezes, a \textbf{reprodução} de respostas corriqueiras sem que haja uma pertinência da proposta e, principalmente, que a solução atenda à dificuldade no âmbito de sua especificidade material.
Em alguns casos, as soluções eram inclusive paradoxais, tendo em vista que a própria solução impunha outro problema para a sustentabilidade.
Como exemplo é possível citar o caso da identificação da ausência de fibras biodegradáveis na produção, cuja solução apontava o uso de fibras 100 \% sintéticas, como aquelas oriundas da reciclagem das garrafas PET.
Os alunos em suas respostas não abordaram algumas dificuldades que a indústria de confecção enfrenta para ser sustentável, por exemplo, a dificuldade da cadeia de produção muito fragmentada.
Muitas respostas afirmaram que o produto sustentável era mal interpretado pelo consumidor, pois era feito de matéria-prima suja.
Acredita -se que o sentido de sujo se relacionasse à ideia de uma matéria-prima já utilizada e reaproveitada.
Outras respostas citaram o pós-consumo, principalmente pelo descarte de roupas compradas ou pelo alto índice de renovação e compras, porém quase nenhuma delas mencionou explicitamente o conceito de economia circular como uma possibilidade de solução.
Algumas vezes, foi encontrada a citação dos brechós como alternativas ao problema, o que representa uma referência implícita à economia circular quando a resposta desenvolve a explicação.
A customização de roupas usadas foi a proposta mais recorrente e que se associava à definição do upcycling, embora este último termo também não tenha aparecido com frequência.
Por meio das respostas, foi possível perceber que o desempenho dos alunos fornece indícios para o pensamento sobre o tema da sustentabilidade no âmbito da formação em Design de Moda.
Pelo nível elementar do desenvolvimento das respostas, verificou -se que a sustentabilidade ainda corresponde a um tema que não parece se fixar como forma de pensamento e, inclusive, cria inconsistências em relação às metodologias projetuais empregadas dentro dos cursos.
A recorrência de respostas superficiais e ainda muito incipientes sobre os produtos do vestuário indicou um distanciamento entre os projetos em design de moda, a formação dos estudantes e a realidade constituinte da indústria do vestuário em nossa sociedade.
A incidência de respostas que indicou a criação de novos produtos a partir de ações muito restritas, como o reaproveitamento de sobras de materiais, pareceu reproduzir discursos que não colocam em pauta a profundidade do tema e a forma de como ele se insere no modo de produção pelo qual a moda opera.
Nesse sentido, o perfil crítico, reflexivo e ético que a questão estimulou não se encontrou refletido nas respostas dos estudantes.
Ao contrário, a formação parece privilegiar uma perspectiva panorâmica e, ainda, muito associada ao perfil de projeto de design de moda que se estabelece pela lógica de produção do novo.
Ou seja, abordagens mais circulares de projetos sustentáveis e que entendam o produto na totalidade de seu ciclo de vida carecem de espaço para desenvolvimento na formação dos estudantes.
Outro aspecto que chamou a atenção foi como as respostas refletem a fragmentação da própria indústria e do lugar do designer na medida em que sugerem o direcionamento e as responsabilidades das ações a terceiros.
Essa constatação demonstra a necessidade de uma associação entre teoria e prática, de modo que a exploração de projetos busque introduzir o designer em uma realidade mais concreta de atuação, na qual ele possa propor novos empreendimentos que dialoguem e provoquem o contexto dominante da própria indústria, refletindo sobre os limites do seu próprio trabalho.
Assim, o perfil de atuação e suas competências serão estimulados na medida em que ele será considerado como um agente, ressaltando sua responsabilidade e seu papel dentro da própria sociedade e da área de conhecimento.
Ainda sobre o aspecto da experiência, é necessário destacar que os índices das respostas mais frequentes evidenciam a falta de uma aproximação com a cadeia têxtil e de confecção.
O desempenho dos estudantes parece revelar pouca familiaridade com setores como os de insumos, da produção e da distribuição, demarcando o designer em uma posição privilegiada e isolada no setor.
Dessa forma, a superficialidade das repostas refletiu uma formação voltada para a criação de produtos pautada por fatores muito herméticos, restringindo -se, algumas vezes, a dimensão estética e apartada da realidade social que o debate da sustentabilidade impõe. 6.3.3 Considerações Finais Através da correção das questões discursivas, pode -se considerar que o nível de desempenho dos estudantes foi regular.
O desenvolvimento das respostas não atingiu os níveis de profundidades esperados, evidenciando conhecimentos ainda pouco estabelecidos na formação dos estudantes.
Pela constituição das questões, foi possível observar que houve uma ênfase no perfil crítico e reflexivo associado à competência de desenvolvimento de produtos de acordo com a metodologia projetual de design e de moda.
As três questões privilegiaram conteúdos que pressupunham uma mesma base de conhecimento das etapas e ferramentas de projeto e da forma de empregá -lo a temas específicos, como a ergonomia aplicada à moda inclusiva e a sustentabilidade.
Esse enfoque caracterizou as questões pela maneira como solicitaram, dos estudantes, a avaliação entre problemas e soluções, o que é muito recorrente na prática profissional do designer.
A vinculação entre variáveis, a capacidade de síntese e a relação de ordenação de etapas examinaram como a formação em design de moda tem abordado algumas capacidades que se associam à dimensão do perfil crítico e reflexivo, mas também criativo e inovador.
Sobre esse aspecto, pode -se afirmar que há algumas carências na formação, principalmente, aquelas associadas às experiências de projetos e à abordagem mais concreta do trabalho do designer.
A insuficiência das respostas evidenciou a falta de uma visão mais profunda da atividade do designer de moda e que se estabelece para além daquela tradicionalmente demarcada pela indústria.
Embora o modelo industrial seja dominante para caracterizar a área do design de moda e, portanto, seus conhecimentos, a abordagem dos temas pelas questões demonstram que é necessário estimular a posição do designer como agente responsável em seu campo de atuação, e que ele está envolvido e conhece as diferentes instâncias que formam a cadeia têxtil e de confecção.
Frente a isso, a \textbf{reprodução} de discursos generalizantes nas respostas põe em pauta um debate sobre a função do designer e, em especial, sua restrição de trabalho na estetização de produtos.
O domínio dos conhecimentos que potencializem a condição interdisciplinar do design e que permitam um reposicionamento de seu lugar no campo de atuação surge como possibilidade de futuros enfoques para área, tendo em vista a valorização de uma formação mais próxima à realidade social contemporânea e aos diferentes modelos de organização do mercado de moda.
Por fim, é necessário pontuar que as respostas dos alunos às questões também sugerem a necessidade de alinhamentos sobre conceitos e conteúdos da área.
A diversidade encontrada nas respostas registra uma condição de precariedade sobre a difusão do conhecimento.
Isso se relaciona à maturidade do desenvolvimento científico da área, ao acesso à informação e ao modo como os cursos desempenham função de produção e \textbf{reprodução} de conhecimento, uma vez que são responsáveis pelo ensino, pesquisa e extensão.
Com isso, deve -se buscar ampliar os espaços de debates e trocas sobre as diretrizes de formação da área, da mesma forma que se deve estimular o aperfeiçoamento constante de todos os envolvidos no ensino.
Por meio dessas ações, talvez se possa chegar a índices mais equânimes sobre a formação dos estudantes, pondo em destaque uma educação atenta ao contexto sociocultural do design de moda.

% matched lemmas: avaliativo, crescer, parâmetro, pedagógico, reprodução
\end{textsample}
